%% [CHANGE: claude-code | 2026-08-20]
%% 007 First Light -- end-to-end forensic and systems case study.
%% Build: pdflatex main && pdflatex main   (twice, for cross-references)
\documentclass[conference,10pt]{IEEEtran}

\usepackage{cite}
\usepackage{amsmath,amssymb}
\usepackage{graphicx}
\usepackage{textcomp}
\usepackage{xcolor}
\usepackage{array}
\usepackage{booktabs}
\usepackage{tabularx}
\usepackage{multirow}
\usepackage{listings}
\usepackage{url}
\usepackage{tikz}
\usetikzlibrary{arrows.meta,positioning,shapes.geometric,fit,backgrounds}

\usepackage[hidelinks]{hyperref}

%% ---- listings style -------------------------------------------------------
\lstdefinestyle{sh}{
  basicstyle=\ttfamily\scriptsize,
  breaklines=true,
  breakatwhitespace=false,
  columns=fullflexible,
  keepspaces=true,
  showstringspaces=false,
  frame=single,
  framesep=3pt,
  rulecolor=\color{black!35},
  backgroundcolor=\color{black!3},
  aboveskip=6pt,
  belowskip=4pt,
  literate={~}{{\textasciitilde}}1
}
\lstset{style=sh}

%% ---- helpers --------------------------------------------------------------
\newcolumntype{L}[1]{>{\raggedright\arraybackslash}p{#1}}
\newcommand{\bs}{\textbackslash}
\newcommand{\tld}{\textasciitilde}
\newcommand{\assertn}[1]{\texttt{assert=#1}}

\begin{document}

\title{Six Opaque Blobs to a Rendered Frame:\\
An End-to-End Forensic and Systems Case Study in Executing an\\
Unmodified Windows AAA Title on Linux Across Two GPU Vendors}

\author{
\IEEEauthorblockN{Shawn}
\IEEEauthorblockA{Luminos OS Project\\
Independent Systems Research\\
Arch Linux / ASUS ROG Zephyrus G14 (GA403UU)}
}

\maketitle

\begin{abstract}
We present a complete, measurement-driven account of taking a 37\,GB set of six
opaque binary containers of unknown provenance and ending with a modern
Direct3D~12 title rendering at native resolution on Linux, on both an integrated
AMD GPU and a discrete NVIDIA GPU. The work spans four normally disjoint
disciplines: binary forensics (establishing that the payload was neither
malicious nor protected by a commercial anti-tamper system), archive
reverse-engineering (extracting and reassembling a container format whose own
installer refused to run), Windows-on-Linux compatibility engineering (a
namespace-fallback defect in Wine, a fixed-path configuration lookup, and a
divide-by-zero in an Optimus display path), and graphics-driver fault isolation
(a segmentation fault inside a proprietary SPIR-V compiler, reached through an
assertion that surfaced as nine unpaintable modal dialogs).

Four findings generalise beyond the specific title. First, code-signing
verification is a \emph{sufficient} argument about anti-tamper presence: if a
retail Authenticode signature still validates over the current bytes, no
protection can have been removed, so the protection question is answered without
any dynamic analysis. Second, a D3D12-to-Vulkan translation layer does not emit
one program per shader; on this workload 4047 of 4430 modules had identical
content hashes but different SPIR-V bytes across vendors, with an NVIDIA-only
extension present in 4039 modules and in zero of the working card's modules.
Third, an entire class of ``hang'' is neither a hang nor a crash but a modal
dialog that no compositor ever paints, and it is diagnosable by counting Win32
window classes rather than by reading logs. Fourth, and most costly: two days of
sophisticated elimination---driver upgrades, cache purges, six Vulkan extensions
disabled, a full shader dump-and-diff, eight code-generation flags---all returned
the identical failure, while the actual fix was to notice that exactly one
translation layer had ever been installed and to substitute a 5\,MB library.
We report the resulting failure taxonomy, the negative results in full, and an
argument for treating unexamined constants as the first bisection axis rather
than the last.
\end{abstract}

\begin{IEEEkeywords}
binary forensics, code signing, entropy analysis, Wine, Proton, Direct3D~12,
Vulkan, SPIR-V, vkd3d, GPU drivers, hybrid graphics, differential diagnosis,
software preservation, silent failure
\end{IEEEkeywords}

\section{Introduction}

\subsection{The problem, stated plainly}

A set of six binary containers totalling 37\,GB arrived on a workstation from a
third party's machine. They were said to contain a 2025-era AAA video game. Three
questions had to be answered, in order, before a single frame could be rendered:

\begin{enumerate}
\item \textbf{Is the payload hostile?} The provenance was a friend's Windows PC,
      and the working hypothesis on arrival was that the files had come from an
      impersonation site rather than the genuine distributor. If any component
      were malicious, the correct action was deletion, not analysis.
\item \textbf{Is the payload protected?} Modern commercial anti-tamper systems
      insert a virtualising layer between the operating system and the program's
      own control flow. Such a layer, or the countermeasures against it, can be
      categorically incompatible with a compatibility layer such as
      Wine~\cite{wine}, in which case no amount of engineering effort would ever
      produce a running program.
\item \textbf{Can it be made to run?} The containers refused to open under their
      own installer, and the title targets Direct3D~12 with Shader Model 6.8,
      ray tracing, and enhanced barriers---the upper bound of what a translation
      layer can currently express.
\end{enumerate}

Each question turned out to have an answer that was cheaper and more decisive
than expected, and each was initially attacked with an unnecessarily expensive
method. That gap---between the method first reached for and the method that
actually settled the question---is the recurring subject of this paper.

\subsection{Why this is worth writing down}

Case studies of this shape are usually either forensic (``here is the malware
we found'') or compatibility-focused (``here is the Wine bug we fixed''). They
are rarely both, and the interesting results here live precisely at the seams:

\begin{itemize}
\item The malware question and the anti-tamper question were \emph{settled by the
      same single measurement}---an Authenticode signature validation---because a
      valid retail signature simultaneously proves nobody added code and proves
      nobody removed any (Section~\ref{sec:denuvo}).
\item The extraction failure that consumed two working sessions was not a
      compression, codec, or archive-integrity problem at all. It was a
      configuration file that a library looked for at one hard-coded absolute
      path, and the ``unsupported compression method'' error it produced was
      technically accurate and completely misleading (Section~\ref{sec:extract}).
\item The final graphics defect presented as neither a crash nor a hang.
      Nine threads were blocked inside modal dialogs that the window system
      created and never painted. The program was, in a strict sense, waiting for
      the user to click a button that had never been drawn on any display
      (Section~\ref{sec:deadlock}).
\item The fix, once found, was to replace a single 5\,MB shared library. It was
      found only after everything sophisticated had failed, because the component
      in question had been an unexamined constant throughout
      (Section~\ref{sec:fix}).
\end{itemize}

\subsection{Contributions}

\begin{enumerate}
\item \textbf{A signature-first argument for anti-tamper triage.} We show that
      Authenticode validation over current bytes is a sufficient and inexpensive
      proof of anti-tamper absence, and that entropy profiling, section-table
      analysis, and import-table analysis serve as independent corroboration
      rather than as the primary instrument (Section~\ref{sec:denuvo}).
\item \textbf{A characterisation of vendor-divergent shader translation.} We
      report, to our knowledge for the first time in the literature, a
      quantitative diff of the SPIR-V emitted by a single D3D12-to-Vulkan
      translator for the same title on two GPU vendors: 4047 of 4430 modules
      share a content hash but differ in bytes, and one vendor-specific SPIR-V
      extension appears in 4039 modules on one card and none on the other
      (Section~\ref{sec:shaders}).
\item \textbf{A diagnostic technique for invisible modal deadlock.} Counting
      windows of Win32 class \verb|#32770| through the compatibility layer's own
      debugger converts an unfalsifiable ``it hangs'' into a scalar that a
      bisection harness can optimise (Sections~\ref{sec:deadlock},
      \ref{sec:harness}).
\item \textbf{A full negative-results table.} Twenty-three distinct hypotheses
      were tested and eliminated by measurement. We publish all of them, because
      the value of this work to a future reader is mostly in what they should
      \emph{not} spend an evening on (Sections~\ref{sec:shaders},
      \ref{sec:fix}).
\item \textbf{A silent-failure taxonomy} drawn from the eleven distinct instances
      encountered here in which a tool reported success, or reported an accurate
      but misdirecting error, while doing nothing useful
      (Section~\ref{sec:lessons}).
\end{enumerate}

\subsection{Scope, ethics, and what was not done}
\label{sec:ethics}

This work was performed on hardware owned by the author, on files already in the
author's possession, for two purposes: determining whether those files were
dangerous, and achieving interoperability with an operating system the vendor
does not support.

Three boundaries are worth stating explicitly, because a reader could otherwise
infer more than was done.

\textbf{No protection was circumvented by this work.} The game executable was
never modified; the central forensic finding is precisely that it is
\emph{byte-identical} to the signed retail build (Section~\ref{sec:denuvo}). The
Steamworks emulation component present in the archive was pre-existing, was
authored by a third party, and is \emph{characterised} here---its file sizes,
its loader hook, its configuration---but was neither written, improved, nor
analysed for the purpose of extending it. Nothing in this paper constitutes a
method for defeating a protection system.

\textbf{The contrasting title was abandoned, not attacked.} A second title on the
same machine does carry a commercial anti-tamper system, and its accompanying
countermeasure is a hypervisor that loads beneath the Windows kernel. We report
this only to explain why that title is categorically out of scope for a
compatibility layer: Wine implements no Windows kernel into which a driver can
load, and on Linux the hardware virtualisation extensions are already owned by
the host. That is an architectural impossibility argument, and it terminated the
investigation rather than directing it.

\textbf{The engineering contribution is entirely on the Linux side.} Every fix
described in Sections~\ref{sec:extract} through \ref{sec:fix}---the namespace
patch, the configuration path, the virtual desktop, the EGL vendor list, the
translation-layer substitution---modifies open-source components or their
configuration. None of them touches the game.

\subsection{Roadmap}

Section~\ref{sec:system} describes the hardware and software under test, which
matters more than usual here because a hybrid-graphics laptop is the source of
at least three of the defects. Sections~\ref{sec:malware} and \ref{sec:denuvo}
answer the provenance and protection questions. Section~\ref{sec:extract} covers
the two extraction failures and the standalone harness built to isolate them.
Section~\ref{sec:reassembly} describes reconstructing a 54\,GB installation
without its installer, and the a-priori verification technique that made every
step checkable in advance. Sections~\ref{sec:graphics1} through \ref{sec:fix}
are the graphics bring-up: reaching a swapchain, a three-session misdiagnosis of
a ``dead'' Vulkan driver, the invisible modal deadlock, the vendor shader
divergence, and the fix. Section~\ref{sec:lessons} generalises.
Section~\ref{sec:appendix} is an artifact and reproduction index.

Throughout, a claim is only stated as established if a specific measurement
supports it, and retracted claims are marked as such rather than deleted---three
of the working hypotheses in this investigation were confidently wrong, and the
manner of their failure is part of the result.

\section{System Under Test}
\label{sec:system}

The platform is load-bearing for this paper. A hybrid-graphics laptop with a
high-DPI panel under a scaling Wayland compositor supplies, by itself, the root
causes of a divide-by-zero in the display path (Section~\ref{sec:graphics1}), a
three-session misdiagnosis of the discrete GPU (Section~\ref{sec:egl}), a
window that appears absent from screenshots while existing (Section~\ref{sec:harness}),
and a game that renders in a box covering 40\% of the panel
(Section~\ref{sec:fix}). None of those would occur on a desktop with one GPU.

\subsection{Hardware}

\begin{table}[htbp]
\caption{Hardware under test}
\label{tab:hw}
\centering
\footnotesize
\begin{tabularx}{\columnwidth}{@{}lX@{}}
\toprule
\textbf{Component} & \textbf{Detail} \\
\midrule
Chassis & ASUS ROG Zephyrus G14, model GA403UU \\
Integrated GPU & AMD Radeon 780M (RDNA~3), no dedicated VRAM; borrows from
                 system memory via the GTT \\
Discrete GPU & NVIDIA GeForce RTX 4050 Laptop, 6141\,MiB dedicated VRAM,
               PCIe-attached, \emph{no attached display} \\
Graphics topology & NVIDIA Optimus (hybrid); the discrete GPU renders and
                    copies, it never scans out \\
System memory & 15.6\,GiB usable, shared bandwidth 102.4\,GB/s with the iGPU \\
Panel & 2880$\times$1800 physical \\
Storage & NVMe; game target on an NTFS volume with 699\,GB free \\
\bottomrule
\end{tabularx}
\end{table}

Two hardware facts recur throughout and are worth isolating now.

\textbf{The discrete GPU has no monitor attached.} On an Optimus system the panel
is wired to the integrated GPU. The discrete GPU is a compute-and-render device
whose output is blitted across PCIe. Any code path that enumerates display modes
for the rendering adapter therefore encounters an empty set, and code that
divides by a refresh rate obtained from that set faults
(Section~\ref{sec:graphics1}).

\textbf{The discrete GPU's memory is dedicated and non-elastic.} Unlike the
Radeon 780M, which borrows from system RAM and can therefore be over-committed
gracefully, the RTX 4050's 6141\,MiB is a hard budget. Anything already resident
on the card is subtracted from what the game can have. During this work a local
language-model service held 3576\,MiB of it, which is 58\% of the card. This did
not cause the central defect---we tested that explicitly and report the negative
result in Table~\ref{tab:eliminated-early}---but it shapes the final launcher
design in Section~\ref{sec:fix}.

\subsection{Software}

\begin{table}[htbp]
\caption{Software stack}
\label{tab:sw}
\centering
\footnotesize
\begin{tabularx}{\columnwidth}{@{}lX@{}}
\toprule
\textbf{Layer} & \textbf{Version / detail} \\
\midrule
OS & Arch Linux, rolling \\
Display & KDE Plasma on Wayland, KWin; scale factor 2 \\
X11 clients & XWayland, \emph{unscaled} \\
NVIDIA driver & 595.71.05, upgraded to 610.57.04 during the work \\
Mesa & 26.1.6 (RADV for the Radeon 780M) \\
Compatibility layer & GE-Proton~10-34; system Wine 11.14 for contrast \\
D3D12 translation & vkd3d-proton (bundled), 2.14.1, and 3.0.1 \\
D3D11/10/9 translation & DXVK \\
Vulkan loader & ICD selection via \verb|VK_ICD_FILENAMES| \\
EGL dispatch & libglvnd, vendor list pinned system-wide \\
\bottomrule
\end{tabularx}
\end{table}

\subsection{The scaling discontinuity}

KDE runs at scale factor 2, so the 2880$\times$1800 panel presents as
1440$\times$900 logical pixels to Wayland-native clients. \textbf{XWayland is
unscaled}, so X11 clients---which includes everything under Wine---work in
physical pixels and see a 2880$\times$1800 display.

This produces two distinct and independently confusing artefacts:

\begin{enumerate}
\item \textbf{Screenshots and compositor coordinates are different number
      systems.} A screenshot is 2880$\times$1800; the compositor's own cursor and
      geometry queries answer in 1440$\times$900. Every pixel measured off a
      screenshot must be halved before it can be used to aim anything. Nothing
      warns about this; the symptom is aiming consistently twice as far right and
      down as intended.
\item \textbf{A correct window looks half-size in compositor reports.} During the
      deadlock investigation, KWin described the game window as
      \verb|960x568|, which reads as a scaling bug. It is correct: a
      1920$\times$1080 physical window is 960$\times$540 logical, plus a 28-pixel
      titlebar. The window was the right size the entire time.
\end{enumerate}

\subsection{The discrete-GPU access gate}

The system enforces a default-deny policy on the discrete GPU. The device node
\verb|/dev/nvidia0| is owned \verb|root:dgpu| with mode 0660, and the
\verb|dgpu| group is deliberately empty. Access is granted only through a setgid
helper, \verb|dgpu-exec-v2|, mode \verb|-rwxr-sr-x root dgpu|. The policy exists
so that the discrete GPU can reach its runtime power-management sleep state
instead of being woken by incidental enumeration.

Two consequences matter for the investigation.

First, \textbf{the gate is the intuitive suspect for every discrete-GPU failure
and was never once the cause}. We record this because it was the first
hypothesis on three separate occasions. It is falsifiable cheaply and in a way
worth reusing: the failures reproduce identically \emph{as root} and with the
gate open, which removes permissions from the space of explanations. In the
central defect the symptom shape also rules it out on its own---a gate failure
produces zero dialogs and a silent fallback to the integrated GPU, whereas
reaching a Vulkan pipeline-creation call at all proves the device was already
opened successfully.

Second, \textbf{a setgid process is non-dumpable}. Anything launched through the
helper has \verb|/proc/PID/mem|, \verb|/proc/PID/maps|, and
\verb|/proc/PID/environ| readable only by root, even for the invoking user.
\verb|/proc/PID/status| remains readable, which is what made the gate's operation
verifiable at all: the game's status file reported
\verb|Gid: 948 948 948 948|, confirming the setgid transition had taken effect.

A related defect found during this work is worth one line, because it is a
general hazard in privilege-raising wrappers: an earlier version of the helper
raised only the \emph{effective} group ID, and \verb|bash| drops that on startup.
Any wrapper invoked through a shell script therefore silently lost the privilege
it had just acquired. The launcher must invoke the v2 helper, which raises the
real GID as well, and it must exec the target binary rather than a shell.

\subsection{Nomenclature used throughout}

\begin{table}[htbp]
\caption{Terms and abbreviations used in this paper}
\label{tab:terms}
\centering
\footnotesize
\begin{tabularx}{\columnwidth}{@{}lX@{}}
\toprule
\textbf{Term} & \textbf{Meaning here} \\
\midrule
\emph{box} & one of the six \verb|.bin| containers \\
\emph{prefix} & a Wine or Proton emulated Windows filesystem and registry \\
\emph{ICD} & Installable Client Driver; a vendor's Vulkan implementation \\
DXIL & the shader bytecode format Direct3D~12 consumes \\
SPIR-V & the shader bytecode format Vulkan consumes \\
\assertn{n} & count of Win32 dialog windows found in the prefix; the
              failure metric defined in Section~\ref{sec:harness} \\
\verb|#32770| & the Win32 window class of a dialog box \\
\bottomrule
\end{tabularx}
\end{table}

\section{Provenance: Is the Payload Hostile?}
\label{sec:malware}

\subsection{The hypothesis on arrival}

The files came from a third party's Windows machine. The stated concern was a
virus, and the specific theory was that the third party had downloaded from an
impersonation site---a lookalike domain serving a reskinned payload under the
name of a known repacking group. That theory has a specific and testable
consequence: an impersonated repack would not match the genuine distributor's
published checksums, and would likely contain persistence mechanisms absent from
the genuine article.

A secondary theory, more sophisticated and more alarming, was raised early: that
the repacking process might involve \emph{re-implementing} a commercial
anti-tamper system, and that such a re-implementation would be an ideal vehicle
for injecting code at runtime. This theory is the reason the anti-tamper question
(Section~\ref{sec:denuvo}) had to be answered even though it initially looked
like a pure compatibility concern.

\subsection{Method}

Analysis was static and offline. The full triage comprised six independent
lines of evidence.

\subsubsection{Checksum provenance}
All six containers were verified against the distributor's published MD5
manifest. All six matched. This does not establish that the distributor is
benign---it establishes that the files are the ones the distributor published,
which directly falsifies the impersonation-site theory.

\subsubsection{Installer script recovery}
The installer is an Inno Setup~\cite{innosetup} package, whose compiled script is
not visible to \verb|strings| and is not the first compressed object in the file.
The recovery procedure is reusable and is documented here because the obvious
approach fails:

\begin{quote}
The \verb|zlb\x1a| magic near the start of the overlay marks the \emph{file
data}, not the script. Locate the SetupLdr offset table by searching for the
12-byte identifier \verb|rDlPtS|, then read nine little-endian 32-bit words:
version, total size, EXE offset, uncompressed EXE size, EXE CRC,
\textbf{offset0} (script header), offset1 (file data), and a table CRC. At
offset0, skip the 64-byte \verb|Inno Setup Setup Data (x.y.z) (u)| banner, then
read a 4-byte CRC, a 4-byte stored size, and a 1-byte compression flag, followed
by chunks of \verb|[crc32(4) + up to 4096 bytes]| until the stored size is
satisfied. The concatenated payload is raw LZMA1 whose first five bytes are the
properties block.
\end{quote}

The decompressed script is the authoritative statement of everything the
installer will do to the system, and reading it converts the persistence question
from a dynamic-analysis problem into a text-search problem.

\subsubsection{Persistence audit}
The recovered script contains exactly one registry write:
\verb|AppCompatFlags\Layers| set to \verb|RUNASADMIN|. There are no \verb|Run|
keys, no \verb|Winlogon| modifications, no service installations, and no
scheduled tasks.

An important negative: string searches for \verb|Startup| and \verb|RegSvr32| do
hit, and in an automated report those hits would be flagged. They are matches
against Inno Setup's own built-in English message table---the strings the
installer would display if it were performing such an action---not evidence that
it performs one. \textbf{A string table is not a call site.} This distinction
accounted for the majority of the initial false positives.

\subsubsection{Network endpoint enumeration}
Every URL in the compiled script resolves to one of three categories: Microsoft
redistributables (Visual C++ runtimes, .NET), NVIDIA PhysX, or the distributor's
own domain. The Inno Download Plugin is present and wired up---so the installer
genuinely can fetch---but the only things it is instructed to fetch are those
redistributables.

\subsubsection{The alarming-but-benign component}
Two files, \verb|host.cmd| and \verb|hosts.exe|, are the sort of thing that
correctly triggers alarm: a program that edits the system \verb|hosts| file.
Reading them settles it. They point roughly sixty lookalike domains at a single
address and then \emph{remove} any existing entry that blocks the distributor's
real domain. This is an anti-impersonation measure authored by the distributor.

The finding is doubly useful. It is benign, and its presence is positive evidence
that the repack is the genuine article---which is exactly the question
Section~\ref{sec:malware} opened with. \textbf{The most alarming artefact in the
package turned out to be the strongest disproof of the theory that motivated the
investigation.}

\subsubsection{Differential comparison against a known-clean package}
Twenty-eight leftover temporary directories were found in the Wine prefix, all
dated to a single day, indicating the installer had been run and had failed
roughly nine times. Every executable and library in those directories was
compared byte-for-byte against the payload of an independently verified clean
repack from the same distributor. All matched, including the archive
decompressor, the download plugin, and the hosts utility.

Four files differed, and each difference is explicable: three are per-repack
configuration (one of which contains the author's own installation path, so it
\emph{must} differ), and the fourth is a newer build of a compression tool with
added codec support.

\subsection{An anomaly correctly excluded}

An earlier pass had flagged a \verb|Winlogon| shell value set to an empty
string---a genuine and serious-looking anomaly. Timestamp comparison excluded it:
the modification predates the arrival of these files by twelve days. It is
unrelated.

We record this because the temptation, when investigating a suspected compromise,
is to attribute every anomaly to the incident under investigation. A system of
any age contains unrelated anomalies. \textbf{Establishing that an artefact is
older than the incident is a complete defence, and it is usually one timestamp
away.}

\subsection{Verdict}

\begin{table}[htbp]
\caption{Malware triage: six independent lines, all negative}
\label{tab:malware}
\centering
\footnotesize
\begin{tabularx}{\columnwidth}{@{}L{2.4cm}X@{}}
\toprule
\textbf{Line of evidence} & \textbf{Result} \\
\midrule
Checksum provenance & All six containers match the distributor's published
                      manifest; impersonation theory falsified \\
Installer script & Recovered in full; one registry write, and it is a
                   compatibility flag \\
Persistence & None. No Run keys, services, scheduled tasks, or Winlogon
              modification \\
Network & All endpoints are Microsoft, NVIDIA, or the distributor \\
\verb|hosts| tooling & Anti-impersonation, authored by the distributor;
                        evidence \emph{for} authenticity \\
Byte comparison & Payload identical to an independently verified clean package,
                  modulo per-repack configuration \\
Runtime footprint & Touches only its own directory, system libraries, and its own
                    application-data folder \\
\bottomrule
\end{tabularx}
\end{table}

\textbf{The files are clean.} No malware, no persistence, no unexpected network
behaviour. The investigation is closed and should not be reopened without new
evidence.

The remaining theory---that the repack re-implements an anti-tamper system, and
that the re-implementation is the injection vector---is not addressed by any of
the above. It is addressed, and destroyed, in the next section, by a single
observation that also happens to answer the compatibility question.

\section{Protection Status: Establishing the Absence of Commercial Anti-Tamper}
\label{sec:denuvo}

This section answers the question the reader most likely arrived for: \emph{how
do you determine, from a binary alone, that a commercial anti-tamper system is
not present?} The answer has a strong form and a weak form, and most published
methodology uses only the weak form.

\subsection{Why the question had to be answered first}

Commercial anti-tamper systems of the Denuvo class operate by taking selected
regions of the program's control flow and replacing them with an interpreter for
a private instruction set, generated per build. The practical consequences for
this project were three:

\begin{enumerate}
\item \textbf{Compatibility.} Such a layer interacts closely with operating
      system internals. A compatibility layer that reimplements those internals
      may or may not satisfy it, and the failure mode is not usually a clean
      error message.
\item \textbf{Effort allocation.} If the protection were present, the ceiling on
      this project might be architectural rather than engineering---in which case
      the correct decision is to stop, as was in fact taken for the contrasting
      title discussed in Section~\ref{sec:contrast}.
\item \textbf{Threat model.} The outstanding malware theory from
      Section~\ref{sec:malware} was that the protection had been
      \emph{reimplemented} by a third party, and that this reimplementation was
      the injection vector. Proving no such layer exists retires that theory
      entirely.
\end{enumerate}

The subject of analysis is \verb|Retail\007FirstLight.exe|, 64{,}228{,}744 bytes,
MD5 \verb|3e0c3e10a9fdb15397f5bc50aeb5e09b|, extracted from the fifth container.

\subsection{The strong argument: signature validation}
\label{sec:sig}

Windows executables may carry an Authenticode signature~\cite{authenticode}: a
PKCS\#7 / CMS structure~\cite{rfc5652} embedded in the certificate directory of
the PE file~\cite{pecoff}, containing a cryptographic digest computed over the
file's contents with three regions excluded (the checksum field, the certificate
table entry, and the certificate data itself).

Validation proceeds in two independent steps, and the distinction between them is
the crux of this section:

\begin{enumerate}
\item \textbf{Chain validation} asks whether the signing certificate descends
      from a trusted root. This says who signed it.
\item \textbf{Digest validation} recomputes the Authenticode hash over the
      \emph{file as it exists now} and compares it to the digest stored inside
      the signature blob. This says whether the file has changed since it was
      signed.
\end{enumerate}

For this executable, both hold. The recomputed SHA-256 equals the digest embedded
in the PKCS\#7 structure. The signer is \textbf{IO~INTERACTIVE~A/S}, via
GlobalSign GCC R45 CodeSigning CA 2020, with a signing time of
2026-07-20 15:43:44 UTC.

\subsubsection{Why this settles the question}

The reasoning is short enough to state completely:

\begin{quote}
The signature validates over the current bytes. Therefore the current bytes are
the bytes the publisher signed. Therefore \textbf{nobody has modified this
executable}---not to add anything, and, crucially, \textbf{not to remove
anything either}. Anti-tamper protection of the class in question is applied to
the executable at build time by the publisher and lives inside the executable's
own sections. If it had ever been present in this build, it would still be
present now, and it would be visible. It is not visible. Therefore it was never
applied to this build.
\end{quote}

This is a \emph{sufficient} argument, and it is the strong form. It requires no
dynamic analysis, no debugger, no emulator, and no unpacking. It is also
self-verifying in a way that heuristics are not: a forged conclusion would
require a hash collision against SHA-256.

It carries a second payload for free. The unresolved malware theory---that the
protection was reimplemented and used as an injection vector---requires
modification of the executable. The signature forecloses that. \textbf{One
measurement answers both questions.}

\subsubsection{The scope of the claim, stated honestly}

The argument establishes that \emph{this executable} is unmodified and
unprotected. It does not establish anything about the other 10{,}949 files in the
package; those were covered separately, by the differential comparison in
Section~\ref{sec:malware}. It also does not establish that no protection exists
anywhere in the product---a separate protected loader could in principle exist as
a distinct binary. Section~\ref{sec:corrob} addresses that by enumerating what
the executable actually imports and what else ships beside it.

\subsection{Corroboration: three independent weak arguments}
\label{sec:corrob}

The strong argument is sufficient, but a single-instrument conclusion is a fragile
conclusion. Three further lines were run, each of which would independently have
raised suspicion had the protection been present.

\subsubsection{Entropy profiling}

Compressed, encrypted, and virtualised code is high-entropy; ordinary compiled
code is not. Entropy analysis is the standard first-pass instrument for detecting
packing~\cite{lyda2007}, with the well-documented caveat that it produces both
false positives and false negatives in the general case~\cite{sokpacker}.

Rather than compute a single figure for the file, which averages away exactly the
signal of interest, the \verb|.text| section was profiled in 43 windows of 1\,MiB
each.

\begin{table}[htbp]
\caption{Shannon entropy profile of the game executable}
\label{tab:entropy}
\centering
\footnotesize
\begin{tabularx}{\columnwidth}{@{}L{3.3cm}X@{}}
\toprule
\textbf{Measurement} & \textbf{Value} \\
\midrule
Whole file & 6.354 \\
\verb|.text|, 43 windows of 1\,MiB & range 5.906--6.700, mean 6.505 \\
Windows at or above 7.5 & \textbf{0} \\
Expected for a virtualised region & 7.9 and above \\
\bottomrule
\end{tabularx}
\end{table}

The distribution is tight, unimodal, and entirely consistent with optimised
native x86-64. There is no high-entropy island anywhere in the code section. Had
a virtualising layer been present, it would appear as a contiguous run of windows
well above 7.5, and windowing is precisely what makes it visible: a whole-file
figure of 6.354 could in principle conceal a small high-entropy region, but 43
individually reported windows cannot.

\subsubsection{Section and import table analysis}

\begin{table}[htbp]
\caption{Structural indicators}
\label{tab:pe}
\centering
\footnotesize
\begin{tabularx}{\columnwidth}{@{}L{3.0cm}X@{}}
\toprule
\textbf{Indicator} & \textbf{Observation} \\
\midrule
Sections & \verb|.text|, \verb|.rdata|, \verb|.data|, \verb|.pdata|,
            \verb|_RDATA|, \verb|.rsrc|, \verb|.reloc| --- stock MSVC layout \\
Anomalous sections & none \\
Imports & 52 DLLs, 753 symbols; a normal compiler-emitted table \\
Runtime & ucrt and \verb|VCRUNTIME140| \\
Protection strings & zero hits for the vendor name, VMProtect, or Themida \\
Size profile & the several-hundred-megabyte bloat characteristic of a protected
               build is absent \\
\bottomrule
\end{tabularx}
\end{table}

Protected builds are structurally distinctive. They typically carry an extra
section with a non-standard name, a truncated or stub import table with the real
imports resolved at runtime, and a substantial size increase over the unprotected
build. None of these is present. The import table in particular is decisive in
its own small way: a protected binary usually cannot afford a full static import
table, because the protection needs to mediate library resolution.

\subsubsection{What the DRM actually is}

Having established what is absent, the analysis identified what is present. The
only protection mechanism in the package is Steamworks licensing, and the
countermeasure is a third-party Steam API emulator.

\begin{table}[htbp]
\caption{The actual protection and countermeasure}
\label{tab:rune}
\centering
\footnotesize
\begin{tabularx}{\columnwidth}{@{}L{3.0cm}X@{}}
\toprule
\textbf{File} & \textbf{Identification} \\
\midrule
\verb|steam_api64.rne| & 319{,}584 B --- the genuine publisher-shipped Valve
                         library, renamed \\
\verb|steam_api64.dll| & 1{,}098{,}392 B --- the third-party replacement \\
\verb|steam_emu.ini| & configuration; carries the author's banner, the
                       application ID, and a configuration section \\
\bottomrule
\end{tabularx}
\end{table}

The mechanism, characterised but not extended (Section~\ref{sec:ethics}), is a
library-loading hook that looks for renamed twins of each library. The
architecturally relevant point is that \textbf{it replaces only Valve's own API
surface}---a documented interface between the game and a storefront---and does
not touch the game's code. This is consistent with, and independently confirms,
the signature finding.

\subsection{The contrasting case, and why it terminated}
\label{sec:contrast}

A second title on the same machine provides the control that makes the above
meaningful: it \emph{does} carry the commercial anti-tamper system, still present
and active. Its accompanying countermeasure is not a library replacement but a
\textbf{hypervisor}: a pair of kernel drivers implementing AMD-V and Intel VT-x
support respectively, which load beneath the Windows kernel and answer the
protection's queries from a layer below it. Source for the virtualisation
component ships in the archive.

This approach cannot work on Linux under a compatibility layer, for two
independent reasons, either of which is sufficient:

\begin{enumerate}
\item \textbf{There is no kernel to load into.} Wine implements Windows APIs in
      user space; it does not implement a Windows kernel, so there is no object
      into which a \verb|.sys| driver can be loaded.
\item \textbf{The hardware is already taken.} The CPU's virtualisation extensions
      are owned by the Linux host. A second hypervisor cannot claim them.
\end{enumerate}

We note this explicitly as an \emph{architectural impossibility}, in the narrow
sense argued in Section~\ref{sec:lessons}: not ``this would be difficult'' but
``this specific mechanism has no counterpart on this platform.'' The
investigation of that title stopped there, which is the correct outcome and the
whole value of asking the protection question before the engineering question.

\subsection{Summary of the method}

For a reader wishing to reproduce this determination on another binary, the
procedure in decreasing order of decisiveness:

\begin{enumerate}
\item \textbf{Validate the Authenticode digest over current bytes.} If it
      validates, the file is the publisher's build and the question is answered.
      Do not skip to entropy; this step is both cheaper and stronger.
\item \textbf{Profile entropy in windows}, not whole-file. Report the window
      count above threshold, and state the threshold.
\item \textbf{Enumerate sections and imports.} A full, compiler-shaped import
      table is difficult for a protection layer to preserve.
\item \textbf{Identify what protection \emph{is} present.} ``No Denuvo'' is not
      ``no DRM,'' and conflating them produces a wrong prediction about runtime
      behaviour.
\end{enumerate}

For this title the answer is unambiguous: \textbf{there is no commercial
anti-tamper system, and there never was one in this build.} Nothing was
stripped, because there was nothing to strip. The engineering path was therefore
open, and the remainder of this paper is about walking it.

\section{Extraction: Two Failures With One Symptom}
\label{sec:extract}

\subsection{The containers}

\begin{table}[htbp]
\caption{Archive inventory, verified by the decompressor's own list command}
\label{tab:inventory}
\centering
\footnotesize
\begin{tabular}{@{}lrrrr@{}}
\toprule
\textbf{Box} & \textbf{Bytes} & \textbf{Files} & \textbf{Packed} & \textbf{Unpacked} \\
\midrule
\verb|fg-01| & 24{,}992{,}387{,}781 & 858   & 23{,}834\,MB & 51{,}419\,MB \\
\verb|fg-02| &  6{,}037{,}379{,}414 & 9{,}774 &  5{,}757\,MB &  5{,}833\,MB \\
\verb|fg-03| &  3{,}591{,}566{,}415 & 1     &  3{,}425\,MB &  3{,}750\,MB \\
\verb|fg-04| &  3{,}280{,}408{,}667 & 218   &  3{,}128\,MB &  4{,}011\,MB \\
\verb|fg-05| &    902{,}897{,}841 & 78    &    861\,MB &  1{,}393\,MB \\
\verb|fg-06| &        525{,}112 & 21    &      0\,MB &      2\,MB \\
\midrule
\textbf{Total} & \textbf{$\approx$37\,GB} & \textbf{10{,}950} & & \textbf{$\approx$64.8\,GB} \\
\bottomrule
\end{tabular}
\end{table}

Note the last row of Table~\ref{tab:inventory}. The smallest container, at
525\,KB and roughly one part in seventy thousand of the total, blocked the entire
project for two working sessions. It contains no game data whatsoever.

A reading trap in the decompressor's own output is worth recording: the size
fields reported by its list command are in \emph{megabytes} for the maximum, and
the running counter is a 32-bit signed integer that goes negative past 2\,GB.
Reading the counter rather than the maximum produces nonsense on any modern
archive.

\subsection{Failure 1: 0.0\% forever}

The installer, run under Wine, stopped at 0.0\% and stayed there. No error, no
timeout, no progress. The naive readings---``it is slow'', ``it deadlocked'',
``the GUI is broken''---are all wrong.

\subsubsection{Root cause}

The repack's external-compressor plugins are split across architectures: a 64-bit
helper process does the work, and a 32-bit plugin loaded inside the installer
communicates with it through shared memory. Tracing the relevant calls with
Wine's relay debugging revealed the mismatch:

\begin{itemize}
\item The 64-bit helper creates its mapping as
      \verb|Global\Read_MapFile_CLS-srep00|.
\item The 32-bit plugin opens \verb|Read_MapFile_CLS-srep00| --- \textbf{with no
      prefix}.
\end{itemize}

On Windows this works, and it works for a reason specific to the NT object
manager~\cite{ntobj}. A session's \verb|\Sessions\N\BaseNamedObjects| directory
is a \emph{shadow} directory: a name that misses there falls through to the
global \verb|\BaseNamedObjects|. \textbf{Wine implements no such fallback.} The
unprefixed open therefore misses, and the plugin responds the way a 2000s-era
Windows utility responds to an unexpected error---it raises a modal message box
reading \verb|Can't open read file mapping!| and waits.

That message box is never painted. The installer is not hung; it is waiting for a
click on a button that exists only in the window manager's data structures. This
is the first of two independent occurrences of this failure shape in this
project; the second, in the graphics stack, is Section~\ref{sec:deadlock}.

\subsubsection{Fix}

One byte per helper. Zeroing the leading \verb|G| of the \verb|"Global\"| string
literal truncates the prefix to the empty string, so both sides use the
session-local name and the open succeeds. Backups of the originals were retained.

\subsubsection{Why patching on disk is insufficient}

This is the part that generalises. Patching the helpers where they live has no
effect, because \textbf{the installer extracts its own fresh, unpatched copies of
every helper into a newly created temporary directory on each run}. A
hand-patched copy elsewhere is simply bypassed.

The working approach is a live patcher: a process that polls at 0.2-second
intervals and patches new copies as they appear. On launch it caught eight
unpatched helpers across two separate temporary directories within seconds.

The general lesson is that \emph{a self-extracting installer is a moving target,
and static patching of its payload is only valid if the payload is not
re-materialised at runtime}. The way to detect this class of problem is to notice
that the patch verifiably applied and verifiably had no effect---otherwise an
extremely confusing pair of facts.

\subsection{Failure 2: ``unsupported compression method''}

With the modal deadlock fixed, five of the six containers opened. The sixth---the
525\,KB one---failed differently:

\begin{lstlisting}
[file] work\work\build01.bat
ERROR: unsupported compression method
ERROR: archive data corrupted (decompression fails)
\end{lstlisting}

\subsubsection{Differential diagnosis: distinguishing the two failures}

The two failures present almost identically to a user: the process is alive and
nothing happens. They are distinguished by a single cheap observation, and
recording that observation is what prevented a second week on the first bug.

\begin{table}[htbp]
\caption{Telling the two extraction failures apart}
\label{tab:diff}
\centering
\footnotesize
\begin{tabularx}{\columnwidth}{@{}L{1.9cm}XX@{}}
\toprule
& \textbf{Failure 1 (modal)} & \textbf{Failure 2 (config)} \\
\midrule
CPU & 0\% & one core fully busy \\
Child process & helper alive & \textbf{none} \\
Bytes written & partial & \textbf{zero} \\
Time to fail & never & $\approx$2\,s \\
Root cause & namespace fallback & missing config file \\
\bottomrule
\end{tabularx}
\end{table}

\subsubsection{The wrong hypothesis, and how long it survived}

The error message is a direct statement that a compression method is not
supported, so the hypothesis was that the container used a codec the shipped
decompressor did not implement. Considerable effort went into it: an archive
scanner was written, the tail directory block was located at two byte offsets,
and a brute-force search over LZMA1 parameter space was implemented to recover
the per-file method table out of the compressed directory block.

None of it was the cause. The archive was correct the entire time.

\subsubsection{Root cause}

The decompression library reads its table of external compressors from a
configuration file, and it looks for that file at \textbf{exactly one hard-coded
absolute location}: the root of the system drive, \verb|c:\arc.ini|. Not its own
directory. Not the current working directory. Not the output directory. Not the
archive's directory.

This was established, rather than guessed, by tracing file operations with
\verb|WINEDEBUG=+file| from three different working directories. Every run
produced exactly one probe, always \verb|CreateFileW L"c:\arc.ini"|, and always
status \verb|c0000034| --- object name not found.

The file exists, in the toolchain directory. With the table absent, every
compression method that depends on an external tool resolves to ``unsupported,''
and extraction dies on the first file that uses one.

\subsubsection{Fix}

One line: copy the configuration file to the root of the emulated system drive.
The container that had blocked the project for two sessions opened on the first
attempt afterwards, as did every other container.

\subsection{Three signals that were available the whole time}

Each of the following was observable before any of the wasted effort began, and
each independently points away from the codec hypothesis.

\begin{enumerate}
\item \textbf{The list operation succeeded throughout.} The directory block
      parsed correctly and reported all files. \emph{A working list beside a
      failing extract points at the method table, not at the archive.} Were the
      archive corrupt or the format misunderstood, listing would fail first.
\item \textbf{Zero bytes written, failing in two seconds, is a parse-time
      failure.} A genuine codec or data fault writes a partial output file
      first, because it fails partway through a stream. Checking the destination
      directory---not the log---separates these instantly.
\item \textbf{``Unsupported compression method'' from this class of tool almost
      always means missing configuration, not a missing codec.} The method table
      is data, and data can simply be absent.
\end{enumerate}

\subsection{Dead ends, published so they are not repeated}

\begin{table}[htbp]
\caption{Failure 2: hypotheses tested and eliminated}
\label{tab:deadends}
\centering
\footnotesize
\begin{tabularx}{\columnwidth}{@{}L{2.7cm}X@{}}
\toprule
\textbf{Hypothesis} & \textbf{How it died} \\
\midrule
Corrupt configuration & Diffed against four copies; all identical \\
Destination filesystem & NTFS versus ext4 made no difference \\
Temporary path setting & Changing it made no difference \\
Regression from the Failure 1 patch & Helpers verified still patched \\
Unknown LZMA1 parameters & Brute-force search found nothing wrong, because
                            nothing was wrong \\
A Wine defect & Two different programs failed on the same file in the same way \\
\bottomrule
\end{tabularx}
\end{table}

\subsection{The standalone harness}

To separate ``the archive cannot be read'' from ``the installer cannot read the
archive,'' a minimal harness was built that loads the decompression library
directly and calls its extraction entry point, bypassing the installer framework
entirely.

The build is constrained in an instructive way. The library is 32-bit, and the
system's Wine is a new-WoW64 build with no 32-bit Unix support, so a 32-bit
\emph{Winelib} binary cannot be produced at all. What can be produced is a 32-bit
\emph{PE}: a freestanding executable with no C runtime and no Windows headers,
built with the LLVM toolchain already present.

\begin{lstlisting}
llvm-dlltool -m i386 -k -d kernel32.def -l kernel32.lib
clang --target=i386-pc-windows-msvc -nostdlib \
      -ffreestanding -O1 -c arcx32.c
lld-link /subsystem:console /entry:start \
      /nodefaultlib /machine:x86 arcx32.obj kernel32.lib
\end{lstlisting}

One implementation detail cost an afternoon and is a general hazard when calling
into legacy Windows libraries: \textbf{the progress callback must be declared}
\verb|__stdcall|. Declared \verb|cdecl|, the stack drifts by 16 bytes per
invocation and the program page-faults after a few hundred callbacks---far enough
into the run to look like a data-dependent fault rather than a calling-convention
error.

\section{Reassembly Without the Installer}
\label{sec:reassembly}

Extraction yields roughly 64.8\,GB of intermediate files across 10{,}950 objects.
These are not the game. They are compressed, delta-encoded, and obfuscated
fragments that the installer's own scripts would normally reconstitute. With the
installer unusable for other reasons, the reconstruction had to be performed by
hand.

The result: a working 54\,GB installation, built in \textbf{37 minutes}, without
ever running the installer.

\subsection{The recipe is inside the archive, in plaintext}

The planned approach was to reverse-engineer the reconstruction procedure out of
the carved installer script (Section~\ref{sec:malware}). That was unnecessary.
The 525\,KB container---the one that had blocked everything---contains the
distributor's own worklists as plain text.

\begin{table}[htbp]
\caption{Contents of the smallest container}
\label{tab:fg06}
\centering
\footnotesize
\begin{tabularx}{\columnwidth}{@{}L{2.0cm}Xr@{}}
\toprule
\textbf{File} & \textbf{Function} & \textbf{Lines} \\
\midrule
\verb|fitgirl01.txt| & restore the audio assets, one command per file & 9{,}774 \\
\verb|fitgirl02.txt| & rebuild the first data chunk & 98 \\
\verb|fitgirl03.txt| & rebuild the second data chunk & 335 \\
\verb|w01/w02.bat| & de-obfuscate halves; a wait-for-file handshake & 9 \\
\verb|build01-03.bat| & \textbf{nothing of value} --- fetch an advert list from a
                        URL shortener and delete themselves & 3 \\
\bottomrule
\end{tabularx}
\end{table}

Two observations follow, and the second is the more useful.

First, \textbf{the file that stopped both the installer and the standalone
harness for two sessions is a three-line advertisement fetcher.} It contains no
build logic at all. The project was blocked on an artefact with no function.

Second, \textbf{the carved installer script was still worth having}, but for a
different reason than expected: the worklists give the commands, and the script
gives the \emph{order} and the \emph{working directory} for each stage, which the
worklists omit. Neither source is sufficient alone.

\subsection{Every tool is a renamed public tool}

The toolchain shipped in the archive consists of eleven executables with opaque
names. None required reverse-engineering. \textbf{Running each with no arguments
caused it to print its own banner.}

\begin{table}[htbp]
\caption{Toolchain identification, by running each binary with no arguments}
\label{tab:tools}
\centering
\footnotesize
\begin{tabularx}{\columnwidth}{@{}L{2.4cm}X@{}}
\toprule
\textbf{Ships as} & \textbf{Actually is} \\
\midrule
\verb|fgpack.exe| & OGGRE v0.1.1, an Ogg Vorbis recompressor~\cite{vorbis} \\
\verb|fgpack2.exe|, \verb|fgpack3.exe|, \verb|lz4.exe| &
  LZ4 at level 12~\cite{lz4}; two of the three are byte-identical \\
\verb|x.exe| & xdelta3, VCDIFF delta decoder~\cite{rfc3284} \\
\verb|xx.exe| & a public XOR utility \\
\verb|x5.exe| & HDiffPatch \verb|hpatchz|, older build~\cite{hdiffpatch} \\
\verb|x5n.exe| & HDiffPatch \verb|hpatchz| v4.8.0, with directory patching \\
\verb|fart.exe| & find-and-replace; expands path placeholders \\
\verb|run.exe|, \verb|countdown.exe| & URL fetcher; sleep \\
\bottomrule
\end{tabularx}
\end{table}

The obfuscation here is nominal, not technical. This is worth stating because the
instinct on encountering eleven unknown binaries inside a package one already
suspects of hostility is to disassemble them. \textbf{The cheapest identification
step---execute with no arguments and read the banner---identified all eleven in
under two minutes}, and each identification was then confirmable against the
public tool's documented command-line grammar.

One tool in the set is worth a note for the forensic reader. The XOR utility
exists because the distributor XOR-obfuscates selected payload files, here with a
64-bit key, purely to evade antivirus heuristics on the compressed intermediates.
XOR is symmetric, so those steps are trivially reversible and cost nothing.

\subsection{A-priori verification: the most valuable technique in this project}

Each directory-level patch is an HDiffPatch directory diff, and \verb|hpatchz|
offers an information mode that prints the diff's manifest \emph{before} anything
is applied: how many input files it expects, their exact total size in bytes, and
the exact size of the output it will produce.

This converts a long pipeline from hopeful to checkable. Every stage could be
verified against a number known in advance, so an error would be caught at the
stage that caused it rather than at the end.

\begin{table*}[t]
\caption{A-priori verification: every stage checked against a manifest read before the stage ran}
\label{tab:verify}
\centering
\footnotesize
\begin{tabularx}{\textwidth}{@{}L{3.2cm}L{6.0cm}X c@{}}
\toprule
\textbf{Stage} & \textbf{Expected (from the manifest)} & \textbf{Measured} & \textbf{Match} \\
\midrule
Audio restore &
  9{,}775 input files, 10{,}456{,}549{,}551\,B total &
  9{,}774 restored files at 6{,}523{,}412{,}777\,B, plus one container at
  3{,}933{,}136{,}774\,B & yes \\
Directory patch output &
  2 files, 12{,}091{,}379{,}792\,B &
  6{,}613{,}152{,}048 + 5{,}478{,}227{,}744 & yes \\
Chunk 0 inputs &
  243 files, 19{,}387{,}609{,}117\,B &
  243 files, 19{,}387{,}609{,}117\,B & yes \\
Chunk 1 inputs &
  411 files, 32{,}085{,}486{,}829\,B &
  411 files, 32{,}085{,}486{,}829\,B & yes \\
Final outputs & --- &
  chunk~0 at 20{,}689{,}753{,}541\,B; chunk~1 at 35{,}798{,}270{,}141\,B & --- \\
\bottomrule
\end{tabularx}
\end{table*}

All five checks landed exactly. The first row is the most informative: the
manifest's expectation of 10{,}456{,}549{,}551 bytes across 9{,}775 files
decomposes precisely into 9{,}774 restored audio files plus one separate
container, to the byte. The 243 and 411 input counts likewise decompose exactly
into their constituent file categories with nothing unaccounted for.

\textbf{The general principle: prefer a pipeline whose stages publish their own
preconditions.} Where such a manifest exists, use it as a checksum you already
have rather than validating only the final artefact---which, at 54\,GB and 37
minutes, is an expensive place to discover a mistake made in minute three.

\subsection{The trap that would have silently broken everything}

The extraction script sent all six containers to a single destination directory.
\textbf{That is wrong for two of them.}

\begin{table}[htbp]
\caption{Per-container extraction destinations}
\label{tab:dest}
\centering
\footnotesize
\begin{tabular}{@{}ll@{}}
\toprule
\textbf{Containers} & \textbf{Destination} \\
\midrule
\verb|fg-01|, \verb|fg-04|, \verb|fg-05| & the application root \\
\verb|fg-02|, \verb|fg-03| & the application's runtime subdirectory \\
\verb|fg-06| & the temporary directory (worklists only) \\
\bottomrule
\end{tabular}
\end{table}

Had this gone unnoticed, the 9{,}774 audio files and their companion container
would have landed one directory too high, and the directory patch in stage 3
would have found no inputs. The failure would have appeared at a stage unrelated
to its cause. The destinations are recorded in the carved installer script, which
is the second reason that artefact remained worth having.

\subsection{The procedure}

Ten stages, in order. Every deletion instruction in the original must be honoured:
the directory diffs consume their inputs, and \verb|hpatchz| deletes the old tree
itself when it merges its temporary output back.

\begin{enumerate}
\item De-obfuscate both halves of the payload (XOR, symmetric).
\item Restore 9{,}774 audio assets, one recompressor invocation each.
\item Apply the top-level directory patch, producing two large intermediates.
\item De-obfuscate both intermediates.
\item Run the chunk-0 worklist.
\item Apply a patch \emph{to the chunk-0 patch}---the directory diff is itself
      shipped as a diff.
\item Apply the chunk-0 directory patch, producing the first data archive.
\item[8--10.] Repeat stages 5--7 for chunk 1.
\end{enumerate}

The per-file shape inside the chunk worklists is consistent: LZ4 to rebuild the
compressed stream, xdelta3 to apply the residual, then XOR, then move into place.

\subsection{Performance: process startup dominates, not compression}

\begin{table}[htbp]
\caption{Measured wall-clock time per stage}
\label{tab:timing}
\centering
\footnotesize
\begin{tabular}{@{}lr@{}}
\toprule
\textbf{Stage} & \textbf{Time} \\
\midrule
De-obfuscate 10.5\,GB & 16\,s (675\,MB/s) \\
9{,}774 recompressor invocations, 6-way parallel & 5\,min \\
Top-level directory patch & 27\,s \\
Chunk-0 worklist & 5\,min 35\,s \\
Chunk-0 archive & 35\,s \\
Chunk-1 worklist & 21\,min \\
Chunk-1 archive & 1\,min 40\,s \\
\midrule
\textbf{Total} & \textbf{$\approx$37\,min} \\
\bottomrule
\end{tabular}
\end{table}

The dominant cost in the two worklist stages is \textbf{Wine process startup, not
compression}. Each individual recompressor invocation performs roughly 40\,ms of
actual work; the remainder is the cost of creating a Windows process under a
compatibility layer. Parallelism is therefore the only tuning knob that matters
for those stages, and a six-way parallel invocation reduced 9{,}774 sequential
launches to five minutes.

One parameter was deliberately reduced from the distributor's own commands: a
cache-size flag was lowered from 4\,GB to 1\,GB. On a 15.6\,GiB machine the
larger value risks the userspace out-of-memory killer, and the change cost
nothing measurable.

\subsection{Result}

A 54\,GB installation with the signature-verified original executable, both
reconstructed data archives, the encrypted package definition the engine expects,
and the Steamworks emulation component in place. The source containers were kept
read-only throughout, so any stage remained redoable: recovery from a mistake is
re-extracting the affected container, roughly 16 minutes for all six.

\section{First Graphics Bring-Up: A Null Pointer With a Documented Cause}
\label{sec:graphics1}

The installation is complete and the executable is the publisher's own. What
follows is the second half of the project, which took considerably longer than
the first, and which is entirely about the graphics stack.

\subsection{Establishing that the reconstruction is actually the game}

Before diagnosing rendering, it was worth confirming that the 54\,GB tree is a
coherent game rather than a coherent-looking pile. Four observations settle it,
and all four come from the game's own startup behaviour rather than from our
inspection:

\begin{itemize}
\item The engine locates and reads its encrypted package definition,
      \verb|packagedefinition.txt| at 35{,}284\,bytes, which is the file the
      Glacier engine consults to find every other asset. A wrong or truncated
      copy stops startup immediately.
\item The crash-telemetry component initialises and writes
      \verb|crash_metrics.dat|, 25{,}240\,bytes.
\item NVIDIA's Streamline framework loads seven \verb|sl.*| libraries and
      \emph{verifies its own digital signatures on each of them}, which is a
      third-party integrity check passing on our reconstruction, performed by
      code we did not write and cannot influence.
\item Streamline then allocates its device-side and host-side pools,
      \verb|sl.chi.heapGPU| and \verb|sl.chi.heapCPU|.
\end{itemize}

The game reaches its rendering path. Everything after this point is a graphics
problem, not a packaging problem.

\subsection{The bundled translation layer is too old}

Direct3D~12 is translated to Vulkan by vkd3d-proton~\cite{vkd3d}. The
system Wine ships a much older vkd3d, and its failure is immediate and legible:

\begin{lstlisting}
Enhanced barriers: 0
\end{lstlisting}

The game requires enhanced barriers, a Direct3D~12 feature for expressing
resource state transitions, and refuses to proceed without them. This is a clean,
early, well-signposted failure---the kind one wants.

Replacing the translation layer with vkd3d-proton 2.14.1 and the Direct3D~11
path with DXVK~\cite{dxvk} changes the log to:

\begin{lstlisting}
DX Ultimate supported!
DXR 1.1 support enabled
Enabling support for SM 6.8
\end{lstlisting}

Ray tracing, mesh shaders, and Shader Model 6.8 all negotiate successfully. The
game then crashes.

\subsection{The crash: a null vtable, disassembled}

The fault is a read from address zero, and the three instructions around it say
exactly what happened:

\begin{lstlisting}
140cb497c:  mov  0x8(%rax),%rcx   ; load an interface ptr
140cb4997:  mov  (%rcx),%rax      ; load its vtable  <-- FAULT
140cb49a8:  call *(%rax)          ; call through it
\end{lstlisting}

This is the canonical shape of an unchecked COM interface query. The game asked
for an interface, did not check the return value, stored the null it received,
and called a method on it one basic block later.

\subsubsection{Which interface}

The requested interface is \verb|ID3D12GraphicsCommandList10|, IID
\verb|{7013c015-d161-4b63-a08c-238552dd8acc}|. It is a recent addition to the
Direct3D~12 interface family. vkd3d-proton 2.14.1 does not implement it and
correctly returns \verb|E_NOINTERFACE|; the game ignores that and dereferences
the null.

\subsubsection{Verifying support without running anything}

A useful and very cheap technique: a COM interface's presence in a translation
layer can be checked by searching the binary for the IID in its
\emph{little-endian in-memory byte order}, which is not the order the GUID is
printed in.

\begin{lstlisting}
grep -qaF $'\x15\xc0\x13\x70\x61\xd1\x63\x4b'\
$'\xa0\x8c\x23\x85\x52\xdd\x8a\xcc' d3d12core.dll
\end{lstlisting}

The first three fields byte-swap; the last two do not. Getting this wrong
produces a confident false negative.

A second false-negative source cost real time and is worth stating plainly:
inside a Proton distribution the Direct3D~12 libraries live under
\verb|files/lib/wine/vkd3d-proton/x86_64-windows|. \textbf{That is
\texttt{lib}, not \texttt{lib64}.} A \verb|lib64| path exists and contains other
things, so a search there returns nothing and looks like an answer.

A third: Wine ships its own \verb|d3d12.dll| of roughly 124\,KB, which is a thin
forwarding stub and is \emph{not} vkd3d-proton. Confirming the size distinguishes
them at a glance.

\begin{table}[htbp]
\caption{Installed translation-layer libraries, for size comparison}
\label{tab:dlls}
\centering
\footnotesize
\begin{tabular}{@{}lr@{}}
\toprule
\textbf{Library} & \textbf{Bytes} \\
\midrule
\verb|d3d12.dll| & 159{,}758 \\
\verb|d3d12core.dll| & 5{,}185{,}550 \\
\verb|d3d11.dll| & 6{,}209{,}550 \\
\verb|d3d10core.dll| & 290{,}830 \\
\verb|dxgi.dll| & 4{,}435{,}982 \\
\verb|nvapi64.dll| & 1{,}974{,}286 \\
\verb|nvofapi64.dll| & 1{,}384{,}462 \\
\bottomrule
\end{tabular}
\end{table}

\subsection{An unrelated fault in the same session: divide by zero in the display path}

Independently of the interface problem, the game faulted while enumerating
display modes. The log sequence is self-explanatory once the hardware topology
from Section~\ref{sec:system} is in mind:

\begin{lstlisting}
Found monitors not associated with any adapter,
  using fallback
readMonitorEdidFromKey: Failed to get EDID
  reg key size
\end{lstlisting}

The discrete GPU has no attached display. Enumerating modes for the rendering
adapter therefore yields an empty set, and the engine divides by a refresh rate
taken from that empty set.

The fix is to give the application a display to find: a Wine virtual desktop, at
which point mode enumeration returns a single well-formed entry and the division
is well defined. This costs nothing and is required on any Optimus-class laptop.

\subsection{The evidence that reversed an hour of reasoning}

The log reported the graphics adapter as vendor \verb|0x1002|, device
\verb|0x73df|---an AMD Radeon. On a machine with two GPUs, running a test
intended to use the NVIDIA card, this reads unambiguously as ``the override did
not take.''

It is the opposite. \textbf{DXVK deliberately reports an AMD device while
running on NVIDIA.} Two options, \verb|dxgi.nvapiHack| and
\verb|dxgi.hideNvidiaGpu|, default to enabled and exist because a number of
shipped games take a worse code path when they detect an NVIDIA adapter through
DXGI. Seeing the AMD identifiers is therefore \emph{positive evidence that the
NVIDIA card is in use}.

This is worth generalising. \textbf{A compatibility layer's job is to lie
convincingly, so its self-reports are the worst available evidence about what
hardware is in use.} Ground truth must come from outside the layer---in this
project, from \verb|nvidia-smi --query-compute-apps|, which reports the process
list the kernel driver actually has.

\subsection{Hypotheses eliminated at this stage}

\begin{table}[htbp]
\caption{Eliminated during first bring-up}
\label{tab:elim-g1}
\centering
\footnotesize
\begin{tabularx}{\columnwidth}{@{}L{2.9cm}X@{}}
\toprule
\textbf{Hypothesis} & \textbf{How it died} \\
\midrule
Vendor-gated crash & Spoofing the vendor as Intel (\verb|8086|) and as NVIDIA
                     (\verb|10de|) both crash at the \emph{identical} address \\
AMD helper library at fault & \verb|amd_ags_x64.dll| is a hard static import;
                              removing it yields \verb|c0000135| before any
                              rendering \\
Never reached rendering & The log shows the swapchain created and presented \\
\bottomrule
\end{tabularx}
\end{table}

The identical crash address under three different spoofed vendors is the
strongest of these: a vendor-conditional code path would branch, and branching
code does not fault at one address.

\subsection{Where this leaves the investigation}

Upgrading the translation layer resolves the missing interface, and the virtual
desktop resolves the display-mode fault. The game then stops crashing---and
starts doing something much harder to diagnose, which is the subject of
Section~\ref{sec:deadlock}. Before that, Section~\ref{sec:egl} deals with a
finding that had to be retracted, because for three sessions the discrete GPU was
believed to be incapable of Vulkan at all.

\section{A Retracted Finding: The Discrete GPU Was Never Broken}
\label{sec:egl}

For three working sessions the recorded conclusion was that the discrete GPU
could not do Vulkan on this system at all---that its driver was, in the words
written down at the time, dead. That conclusion was wrong, it was wrong for a
reason that generalises, and correcting it is what made everything in
Sections~\ref{sec:deadlock} through~\ref{sec:fix} possible.

\subsection{The symptom}

Any attempt to create a Vulkan instance restricted to the NVIDIA driver failed
at the earliest possible point:

\begin{lstlisting}
Could not get 'vkCreateInstance' via
 'vk_icdGetInstanceProcAddr' for ICD libGLX_nvidia.so.0
\end{lstlisting}

The loader could open the driver library but could not obtain a single entry
point from it. Every global function resolved to null. This is not a subtle
failure; it looks exactly like a broken or mismatched driver.

\subsection{Isolating it below the compatibility layer}

The first useful decision was to stop testing through Wine. A twenty-line C
program was written that does nothing but \verb|dlopen| an ICD library, call its
negotiation entry point across every interface version the loader supports, and
attempt one instance creation.

\begin{table}[htbp]
\caption{Direct ICD probe, outside Wine and outside the Vulkan loader}
\label{tab:icd}
\centering
\footnotesize
\begin{tabularx}{\columnwidth}{@{}L{3.4cm}cc@{}}
\toprule
\textbf{Probe} & \textbf{Mesa pin} & \textbf{NVIDIA present} \\
\midrule
\verb|icdNegotiate| v1--7 & $-3$ & $0$ \\
\verb|vkCreateInstance| & $-3$ & $0$ \\
Resulting instance handle & \verb|(nil)| & \verb|0x557e90579430| \\
\bottomrule
\end{tabularx}
\end{table}

The difference between the two columns is a single environment variable. That is
the whole finding: the driver is fine, and something in the environment is
disabling it.

\subsection{Root cause: an environment variable that replaces rather than adds}

The system pins its EGL vendor selection through
\verb|__EGL_VENDOR_LIBRARY_FILENAMES|, set to Mesa's \verb|50_mesa.json| so that
ordinary desktop applications do not incidentally wake the discrete GPU.

\textbf{That variable is a replacement, not an addition.} It does not extend the
vendor search path; it \emph{is} the vendor search path. With only Mesa's entry
visible, libglvnd~\cite{glvnd} has no NVIDIA vendor at all.

The non-obvious part is why this breaks \emph{Vulkan}, which is a separate API
with its own loader and its own ICD list. NVIDIA's Vulkan implementation
initialises through the same vendor-dispatch machinery as its EGL
implementation. With the vendor absent, that initialisation fails silently and
the ICD returns a table of null pointers rather than an error.

The confirming observation is from \verb|strace|: in the failing configuration
the driver library \textbf{never opens \texttt{/dev/nvidiactl}}. It is not being
refused by the kernel; it never asks. This also independently eliminates the
access gate of Section~\ref{sec:system}, which was again the first suspect.

\subsection{Fix}

Whenever the NVIDIA Vulkan driver is required, the EGL vendor list must name
\emph{both} vendors---Mesa's entry and \verb|60_nvidia.json|. The rule to carry
forward:

\begin{quote}
\textbf{A pin that names one vendor is a policy statement, and every consumer of
that policy inherits it, including ones in a different API.} Any launcher that
needs the pinned-out vendor must restore it explicitly rather than assume it can
add to a list.
\end{quote}

After the correction, \verb|vulkaninfo| reports the discrete GPU exactly as it
should: \textbf{NVIDIA GeForce RTX 4050 Laptop GPU}, device type
\verb|DISCRETE_GPU|, Vulkan 1.4.329, driver 595.71.05.

\subsection{A second, independent fault hiding behind the first}

Correcting the vendor list revealed that a second defect had been masked by it.
The project's Wine launcher, when asked for the discrete GPU, was performing a
bare \verb|exec wine| with no ICD selection at all---so it silently used the
integrated Radeon while the desktop notification it emitted said NVIDIA.

Two independent faults, one of which lies to the user about which hardware is in
use, are a good illustration of why the ground-truth rule from
Section~\ref{sec:graphics1} matters. The failing configuration was
indistinguishable from the working one by every self-report available, and
distinguishable instantly by \verb|nvidia-smi --query-compute-apps|.

\subsection{The methodological lesson}

The recorded claim was ``NVIDIA Vulkan is dead on this machine.'' The evidence
for it was a test that had been run under a system-wide policy the test did not
know about.

\begin{quote}
\textbf{When an instrument reports that a major component is entirely
non-functional, suspect the instrument first.} Total failure is rare; misconfigured
measurement is common. The cheapest discriminator is to reproduce the failure at
the lowest possible layer---here, forty lines of C with no loader, no
compatibility layer, and no launcher script between the test and the driver.
\end{quote}

A corollary specific to this class of bug: a bare test of the form ``set the ICD
list and see whether it works'' is \emph{not} a test of the driver, because the
ICD list is not the only list involved. It is a test of one variable in an
environment with at least two.

With the discrete GPU restored to service, the game could finally be run on it.
It then produced the failure that defines the rest of this paper.

\section{The Central Defect: Nine Invisible Dialogs}
\label{sec:deadlock}

This is the failure that consumed most of the project. It is also the second
independent occurrence of the failure shape first met in
Section~\ref{sec:extract}: \emph{a program that is not hung, waiting on a window
nobody can see}.

\subsection{Symptom}

On the discrete GPU, the game launches and appears to work. Specifically:

\begin{itemize}
\item A window opens, correctly sized, and stays open.
\item Roughly 1{,}960\,MiB of video memory is allocated on the discrete GPU.
\item The expected worker threads exist: \verb|vkd3d-swapchain|,
      \verb|vkd3d_queue|, \verb|vkd3d_fence|.
\item Then GPU utilisation falls to \textbf{0\%} and the shader clock drops to
      its idle 210\,MHz.
\item The window is black. It stays black indefinitely.
\item \textbf{Nothing further is written to any log.}
\end{itemize}

The process is alive, holding memory, holding a device, and doing nothing. No
error, no crash, no timeout.

\subsection{Making the invisible visible}
\label{sec:harness}

The breakthrough was not a hypothesis. It was asking the window manager inside
the compatibility layer to enumerate its own windows:

\begin{lstlisting}
winedbg --command "info wnd"
\end{lstlisting}

There were \textbf{nine} windows of class \verb|#32770|---the Win32 dialog box
class---each containing three buttons reading \verb|&Abort|, \verb|&Retry|,
\verb|&Ignore|. Nine assertion dialogs, stacked, none of them ever painted to the
screen, each blocking a thread on a click that will never arrive.

The assertion text is:

\begin{lstlisting}
../src-wine/dlls/winevulkan/loader_thunks.c:3151
!status && "vkCreateComputePipelines"
\end{lstlisting}

\subsubsection{The detail that redirected the entire investigation}

The natural reading is that a Vulkan call returned an error, and the natural next
step is to capture the \verb|VkResult|. That is a category error, and recognising
it saved an unknown amount of further wasted effort.

\textbf{\texttt{status} here is an NTSTATUS describing whether the call into the
Unix-side driver \emph{completed}, not a VkResult describing what it
returned.} Its value is \verb|0xc0000005|: an access violation. The driver did
not reject the pipeline. \textbf{The driver crashed while compiling it}, and the
thunk layer is reporting the crash of the call itself.

Turning on structured-exception tracing gives the address:

\begin{lstlisting}
Exception 0xc0000005 at 0x7f04a4746fd8:
  libnvidia-glvkspirv.so.610.57.04 + 0x346fd8
\end{lstlisting}

Nine threads, nine faults, one address---inside NVIDIA's SPIR-V shader compiler.
This is not a game bug and not a translation-layer bug in the ordinary sense. It
is a driver segmentation fault, reached deterministically.

\subsubsection{A cheaper proxy for the same signal}

Attaching a debugger to every test run is slow. A far cheaper indicator turned
out to be already present in the ordinary log:

\begin{lstlisting}
fixme:oleacc:find_class_data unhandled window
  class: L"#32770"      (x29)
fixme:oleacc:find_class_data unhandled window
  class: L"Button"      (x9)
\end{lstlisting}

The accessibility layer announces every window class it does not recognise, so
counting those lines counts the dialogs without a debugger. This became the
project's primary metric.

Two traps in reading the debugger output are worth recording. \verb|winedbg|
truncates window text to fourteen characters, so button labels appear cut off and
look corrupted. And the game window is reported with a geometry of
\verb|960x568|---which reads as a scaling bug and is not one, for the reason
given in Section~\ref{sec:system}. It was also simply sitting behind a browser
window.

\subsection{The differential that names the culprit}

The same build, the same game, the same command line, on the integrated Radeon:
it renders. On the discrete NVIDIA card: nine faults in the shader compiler.

One variable differs, and it is the driver. Everything after this point is about
\emph{why} the same program produces shader code that one vendor's compiler
survives and the other's does not---which is answered, unexpectedly, in
Section~\ref{sec:shaders}.

\subsection{An automated harness, and why it had to be negative-tested}

Manual testing of a failure that takes 40--90 seconds to manifest, has no error
message, and must be counted rather than observed, is not viable. A harness was
built:

\begin{lstlisting}
007-try.sh <label> [--igpu] [--proton=NAME] \
           [--wait=N] [--kill] [VAR=VAL ...]
\end{lstlisting}

It launches one configuration, waits, samples, kills, and emits a single line:

\begin{lstlisting}
alive=Y assert=9 gpu=0 rss=2107448 secs=62
\end{lstlisting}

\textbf{The harness was then run against a configuration already known to fail,
to confirm it reported failure.} This is not ceremony. An earlier version of the
assertion counter used a construct of the form

\begin{lstlisting}
N=$(grep -c 'pattern' "$LOG" || echo 0)
\end{lstlisting}

which is wrong in a way that is invisible: \verb|grep -c| prints \verb|0| and
\emph{also} exits non-zero when there are no matches, so the fallback appends a
second zero and \verb|N| becomes the two-line string \verb|"0\n0"|. Every
subsequent numeric comparison then errors and evaluates false, and the guard
never fires. A health check that cannot report ill health is worse than no health
check, because it is trusted.

\subsection{Hypotheses eliminated before the shader work began}

\begin{table}[htbp]
\caption{Early eliminations, each with the measurement that killed it}
\label{tab:eliminated-early}
\centering
\footnotesize
\begin{tabularx}{\columnwidth}{@{}L{2.5cm}X@{}}
\toprule
\textbf{Hypothesis} & \textbf{Measurement} \\
\midrule
Still compiling shaders &
  The 259\,MB pipeline cache's file descriptor position reaches the file size
  exactly, then stops. A second run with the cache fully warm reaches the
  identical stall in 46\,s \\
Out of video memory &
  Stopping the resident language-model service freed 3{,}576\,MiB of 6{,}141.
  Identical stall \\
Online check or licensing &
  \textbf{Zero} TCP sockets open in the process at the time of the stall \\
Missing virtual desktop &
  Present and correct; the display-mode fault of
  Section~\ref{sec:graphics1} does not recur \\
Missing Vulkan extension &
  All six extensions the translation layer requests are present and reported by
  the driver \\
Discrete-GPU access gate &
  Reproduces identically as root with the gate open
  (Section~\ref{sec:system}) \\
\bottomrule
\end{tabularx}
\end{table}

The first row deserves emphasis because it is the hypothesis that survives
longest by default. ``It is still working'' is unfalsifiable without a
measurement, and the measurement here is unusually direct: the operating system
will report how far into a file a process has read. When that offset equals the
file size and stops advancing, reading is finished, and ``still loading'' is
dead.

\subsection{A trap in process matching}

The harness matches the game process by name. On Linux the kernel's \verb|comm|
field is truncated to fifteen characters, so the process is
\verb|007FirstLight.e| and never \verb|007FirstLight.exe|. A pattern containing
the full name matches nothing, forever, silently---and the harness reports
\verb|alive=N| for a process that is running perfectly well.

The related hazard, which bit earlier in the project, is that a substring match
against the full command line (\verb|pgrep -f|) matches the shell that is
performing the search. The correct tool is an exact-name match against the
truncated \verb|comm|.

\section{The Shaders: Same Game, Different Bytes Per Vendor}
\label{sec:shaders}

The driver crashes compiling a shader. The obvious next question is \emph{which}
shader, and the obvious assumption is that both GPUs receive the same shaders and
only one vendor's compiler mishandles them. The second half of that assumption is
false, and correcting it reframes the problem.

\subsection{Dump both, and compare}

The translation layer can be told to write every shader it produces, in both its
input and output forms, to a directory. Running the game once on each GPU and
comparing the two directories gives the following.

\begin{table}[htbp]
\caption{Shader output compared across the two GPUs, same game, same build}
\label{tab:shaders}
\centering
\footnotesize
\begin{tabularx}{\columnwidth}{@{}L{4.4cm}X@{}}
\toprule
\textbf{Measurement} & \textbf{Value} \\
\midrule
Modules translated, Radeon & 6{,}657 \\
Modules translated, NVIDIA & 4{,}430 \\
Hashes present in both sets & 4{,}430 \\
Same hash, \textbf{different bytes} & \textbf{4{,}047} \\
Uses \verb|SPV_NV_raw_access_chains| & 4{,}039 NVIDIA / \textbf{0} Radeon \\
\quad of which compute shaders & 561 of 886 \\
Uses \verb|SPV_NV_shader_|\linebreak\verb|subgroup_partitioned| & exactly 1
                                                                 compute shader \\
\bottomrule
\end{tabularx}
\end{table}

The interpretation is unambiguous. \textbf{The translation layer emits
vendor-specific code.} Given identical Direct3D shader bytecode as input, it
produces different Vulkan output depending on which driver is present, because it
uses vendor extensions where they are available.

Of the shaders present on both sides, \textbf{91\% differ byte-for-byte}. The
NVIDIA path uses a vendor extension in 4{,}039 modules that the Radeon path uses
in none. So the premise ``the same shaders work on AMD'' was never true. AMD was
never given these shaders.

This matters for a reason beyond diagnosis: it means the AMD success provided
much weaker evidence than it appeared to. It did not show that the game's shaders
are sound; it showed that a \emph{different} set of shaders is sound.

\subsection{Overriding the shaders, and a large partial result}

The translation layer also supports substituting a shader by hash: place a
replacement in a directory and the corresponding module is loaded from there
instead of being generated. Feeding the Radeon's modules to the NVIDIA run is
therefore mechanically possible.

The result was the single largest improvement of the investigation:

\begin{center}
\texttt{9 dialogs, GPU 0\%} $\longrightarrow$
\textbf{\texttt{4 dialogs, GPU 100\%}}
\end{center}

Five of the nine crashing pipelines stopped crashing, and the discrete GPU went
from idle to fully loaded---it was doing real work for the first time. The result
held across \textbf{fourteen consecutive invariant runs}, so it is not noise.

It is also not a fix. Four dialogs remain, and the configuration is not one that
could ship: it requires a shader dump from the other GPU, and it forces the
translation layer to bypass its own pipeline cache.

\subsection{Why this experiment is not clean, stated plainly}

It is tempting to read the partial success as ``the vendor extension is the
bug.'' The experiment does not support that, for two reasons that must be stated:

\begin{enumerate}
\item Substituting shaders \textbf{forces the layer to ignore its cached SPIR-V
      entirely}, changing a second variable at the same time as the first.
\item The Radeon modules were compiled for a different device's capability set.
      Feeding them to a different vendor's compiler is not a controlled
      substitution; it is a different program.
\end{enumerate}

\subsection{An unresolved contradiction, published rather than smoothed over}

If the vendor extension were the cause, then disabling that extension outright
should reproduce the improvement. It does not.

\begin{quote}
Disabling \verb|VK_NV_raw_access_chains| through the layer's own extension
filter---so that the same 4{,}039 modules are generated without it---does
\textbf{not} reduce the dialog count. The override, which changes both the
extension usage and the cache path, does.
\end{quote}

We were unable to resolve this contradiction, and we report it unresolved. The
most likely explanations are that the improvement came from the cache bypass
rather than the extension, or that the extension filter does not affect
already-cached modules. Neither was confirmed. It is recorded here because a
reader reproducing this work will hit the same contradiction, and because the
correct fix, found later, makes the question moot in a way that is itself
informative: the layer that emitted the shaders was the wrong layer to be
adjusting.

\subsection{Everything tried at the shader level}

\begin{table}[htbp]
\caption{Shader- and translation-layer configuration attempts, all negative}
\label{tab:shaderelim}
\centering
\footnotesize
\begin{tabularx}{\columnwidth}{@{}L{4.2cm}X@{}}
\toprule
\textbf{Attempt} & \textbf{Result} \\
\midrule
Disable the vendor extension & no change \\
Disable subgroup-partitioned ops & no change \\
Force single-threaded pipeline compilation & no change \\
Disable the pipeline library & no change \\
Disable pipeline caching entirely & no change \\
Force synchronous shader compilation & no change \\
Disable descriptor-buffer usage & no change \\
Disable ray tracing & no change \\
Disable mesh shaders & no change \\
Force a conservative memory model & no change \\
Reduce shader-compiler thread count & no change \\
Substitute the other vendor's modules & \textbf{9 $\to$ 4 dialogs, GPU to
                                        100\%} \\
\bottomrule
\end{tabularx}
\end{table}

Eleven configuration changes with no effect and one with a large effect is itself
a signal. When a component's entire exposed configuration surface fails to move
the failure, the component is probably not the one to be configuring.

\subsection{A verification script that was itself wrong}

The override mechanism prints a line for each substituted module, so the harness
checked for those lines to confirm the override had taken effect. It reported
success on runs where nothing had been substituted, because of the \verb|grep -c|
defect described in Section~\ref{sec:deadlock}. Two sweep results had to be
discarded and re-run.

\begin{quote}
\textbf{Every instrument in this project that was not deliberately tested against
a known-bad input was, at some point, wrong.} The harness, the override
verifier, the shader-count comparison, and the process matcher each produced a
confident false reading before being corrected.
\end{quote}

At this point the shader hypothesis was exhausted: a real and reproducible
improvement, no clean explanation, and no path to a shippable configuration. The
next section changes a variable that had been held constant since the beginning
of the project, and the problem disappears entirely.

\section{The Fix: Changing the Variable Nobody Was Varying}
\label{sec:fix}

\subsection{The unexamined constant}

Every experiment in Sections~\ref{sec:deadlock} and~\ref{sec:shaders} varied
something inside the translation layer or the driver, and every one of them held
the \emph{compatibility distribution} fixed. GE-Proton~10-34 had been chosen at
the start because it worked, and it was never afterwards treated as a variable.

A compatibility distribution bundles its own build of the Direct3D~12 translation
layer. Changing the distribution changes that layer. The question ``does a
different build of the translation layer produce shaders the driver survives?''
had therefore never been asked, despite the entire investigation being about
shaders the driver does not survive.

\subsection{Building a variant cheaply}

Replacing the bundled layer in place would destroy a working configuration. The
approach taken instead was to clone the distribution and replace only the two
libraries of interest.

\begin{lstlisting}
cp -al "$SRC" "$DST"                 # hardlink clone
for d in <four library directories>; do
  rm -f "$DST/$d/d3d12.dll" "$DST/$d/d3d12core.dll"
  cp    "$NEW/d3d12.dll"    "$DST/$d/"
  cp    "$NEW/d3d12core.dll" "$DST/$d/"
done
\end{lstlisting}

A hardlink clone of a 1.4\,GB distribution costs about \textbf{25\,MB}, because
only the replaced files are new data. This is safe \emph{only} because the
replacement is \verb|rm| followed by \verb|cp|---creating a new inode---rather
than an in-place edit, which would modify the original through the shared inode.
That distinction is the entire safety argument and is easy to get wrong.

Two further points, each of which cost a wasted result:

\begin{itemize}
\item \textbf{A prefix upgrade is one-way.} Once an emulated Windows filesystem
      has been upgraded by a newer distribution it cannot be used by an older
      one. Each variant needs its own prefix.
\item \textbf{Seeding a prefix requires more than the prefix directory.} Copying
      only \verb|pfx/| omits the sibling \verb|tracked_files| and \verb|version|
      metadata, and the launcher aborts before the game starts---which reads
      exactly like the variant being broken. This defect made the
      \emph{winning} configuration look like a failure on its first run.
\end{itemize}

\subsection{Result}

\begin{table}[htbp]
\caption{Translation-layer builds, all other variables held fixed}
\label{tab:proton}
\centering
\footnotesize
\begin{tabular}{@{}llrr@{}}
\toprule
\textbf{Build} & \textbf{Layer} & \textbf{Dialogs} & \textbf{GPU} \\
\midrule
Stock (as used throughout) & 2.14.1 & 9 & 0\% \\
Clone with newer layer & \textbf{3.0.1} & \textbf{0} & \textbf{100\%} \\
Clone with same layer & 2.14.1 & --- & exits early \\
Newer distribution & 3.x & 0 & 34\% \\
\bottomrule
\end{tabular}
\end{table}

\textbf{Zero dialogs. The game runs.} The driver's shader compiler does not
crash, because it is no longer being handed the code that crashed it.

The third row is the control: the same clone procedure with the \emph{original}
layer does not fix anything, which rules out the cloning itself as the cause of
the improvement.

\begin{table}[htbp]
\caption{Before and after, discrete GPU}
\label{tab:beforeafter}
\centering
\footnotesize
\begin{tabularx}{\columnwidth}{@{}L{3.0cm}XX@{}}
\toprule
& \textbf{Before} & \textbf{After} \\
\midrule
Assertion dialogs & 9 & \textbf{0} \\
GPU utilisation & 0\% & \textbf{100\%} \\
Video memory & 1{,}960\,MiB, static & 2{,}064--2{,}070\,MiB, live \\
Shader clock & 210\,MHz idle & boost \\
Window & black, indefinitely & title screen at 2880$\times$1800 \\
\bottomrule
\end{tabularx}
\end{table}

\subsection{Verified three independent ways}

Following the project rule that a fix is not a fix until it has been exercised
end to end, the result was confirmed by three methods that do not share a failure
mode:

\begin{enumerate}
\item A \textbf{199-second continuous run} with the kernel driver's own compute
      process list sampled throughout---ground truth from outside the
      compatibility layer, per Section~\ref{sec:graphics1}.
\item A \textbf{screenshot of the in-game graphics menu}, with DLSS listed among
      the upscaling options. That option is only offered when an NVIDIA adapter
      is present, so the game itself corroborates which GPU it is on.
\item A \textbf{screenshot of the title screen} taken from the installed
      launcher command as an ordinary user would invoke it---not from the test
      harness.
\end{enumerate}

The third matters more than it looks. A fix that works only under the harness is
a fix to the harness.

\subsection{The integrated GPU was re-tested, not assumed}

The Radeon path was re-run under the new configuration rather than presumed
intact: \verb|assert=0|, 95\% utilisation. Both GPUs work. The launcher defaults
to the discrete card and accepts a flag to select the integrated one.

\subsection{Launcher design decisions}

Three small decisions in the shipped launcher are worth recording because each
encodes something learned above.

\textbf{Restore the EGL vendor list explicitly.} Per Section~\ref{sec:egl}, the
system-wide pin must be widened to include the NVIDIA vendor, not assumed to be
additive.

\textbf{Detect the panel's native mode from the kernel, not from a
configuration file.} A subtlety: files under \verb|/sys| report a size of zero
regardless of content, so the common shell idiom for ``file exists and is
non-empty'' is always false there. The file must be read and its content checked.

\textbf{Report memory pressure; do not act on it.} The discrete card's memory is
a hard 6{,}141\,MiB budget, and a resident language-model service commonly holds
3{,}576\,MiB of it. The launcher prints a notice when free memory is low. It does
\textbf{not} stop the service, because that service may be doing work the user
cares about more than the game. Tools that silently kill other work to make room
for themselves are a category of behaviour to avoid.

\subsection{Closing the remaining lead}

One lead from the original three-way plan remained: try a newer graphics driver.
It is closed negatively---all available prerelease driver packages resolve to the
same version already installed, 610.57.04. There is no newer driver to try. The
lead is exhausted rather than unexplored, which is a different and more useful
state to record.

\subsection{Undo}

The original distribution is untouched; the variant is a separate directory with
its own prefix. Reverting is deleting two directories and one line in the
launcher. The fix is additive, and the cost of being wrong about it is a few
seconds.

\section{Discussion: What Generalises}
\label{sec:lessons}

Nothing in this project required a novel technique. What it required, repeatedly,
was noticing that a question had not been asked. This section extracts the parts
that transfer.

\subsection{Search order dominates search effort}

The defect was in a component that was never varied, and it was found within one
attempt of finally varying it. The preceding weeks were spent varying components
that could not have been responsible.

The failure of process is identifiable in advance. \textbf{A component chosen at
the start of a project because it worked becomes invisible to the project.} It is
not on the list of suspects because it is part of the definition of the setup.

The practical rule follows from delta debugging~\cite{zeller2002}, applied one
level up from where it is usually applied:

\begin{quote}
\textbf{Before bisecting within a component, enumerate the components you have
never bisected across.} Any element of the stack that has held one value for the
whole investigation belongs on the suspect list \emph{ahead of} the elements you
have already varied without success, not behind them.
\end{quote}

A corroborating signal was available and misread: eleven configuration changes
inside the translation layer produced no effect at all (Table~\ref{tab:shaderelim}).
An exhausted configuration surface is evidence that the component is the wrong
one to configure---not an invitation to configure it harder.

\subsection{Differential diagnosis beats hypothesis generation}

Twice in this project, two failures shared a symptom and were separated by one
cheap observation:

\begin{itemize}
\item Extraction: process alive, nothing happening. \emph{Is a core busy, and did
      it write any bytes?} Two seconds of observation, two entirely different
      root causes (Table~\ref{tab:diff}).
\item Rendering: window open, nothing happening. \emph{Ask the window manager
      what windows exist.} One command, and an invisible modal dialog becomes a
      driver crash address.
\end{itemize}

In both cases the correct move was to \emph{enumerate the state} rather than to
generate hypotheses about it. Hypothesis generation is unbounded; state
enumeration is finite and usually one command.

\subsection{Prefer pipelines whose stages publish their preconditions}

The reassembly (Section~\ref{sec:reassembly}) succeeded on the first attempt
across ten stages and 54\,GB because every stage could be checked against a
manifest read \emph{before} it ran. Five checks, five exact matches, and any
error would have surfaced at its own stage.

Where a format offers this, it is the cheapest correctness guarantee available:
a checksum you did not have to compute and did not have to trust yourself to
choose.

\subsection{Instruments lie, and only negative testing catches it}

Every measurement tool built in this project produced at least one confident
false reading before being corrected:

\begin{itemize}
\item The assertion counter reported zero because \verb|grep -c| both prints zero
      and fails.
\item The override verifier reported success on runs that overrode nothing, for
      the same reason.
\item The process matcher reported the game dead because \verb|comm| truncates at
      fifteen characters.
\item The graphics layer reported an AMD adapter while running on NVIDIA, by
      design.
\item The launcher's own notification claimed NVIDIA while using the integrated
      GPU.
\item The recorded conclusion ``NVIDIA Vulkan is non-functional'' was an artefact
      of a system-wide policy the test was unaware of.
\end{itemize}

\textbf{The discipline that catches all six is the same: run every instrument
against an input whose answer you already know, and specifically against a
\emph{failing} input.} A health check that has only ever been shown reporting
health has not been tested.

\subsection{``Impossible'' has exactly one legitimate use}

This paper contains one genuine impossibility claim---the hypervisor-based
countermeasure of Section~\ref{sec:contrast}, which fails for two independent
structural reasons---and it is the only place the word is used.

Elsewhere, several things that were recorded as impossible turned out to be
merely unexamined: a compiled installer script that ``cannot be read'' was read
by parsing the format; a driver that was ``dead'' was disabled by an environment
variable; an archive that used ``an unsupported codec'' was missing a
configuration file.

\begin{quote}
\textbf{Separate three claims that are easily conflated: this cannot exist; this
would be difficult; I do not currently know how.} Only the first is a stopping
condition, and it requires a structural argument, not a failed attempt.
\end{quote}

\subsection{The strong form of a negative result}

Section~\ref{sec:denuvo} makes a claim of absence---that no commercial anti-tamper
system was ever in this build---and supports it with a cryptographic identity
rather than an accumulation of heuristics.

This is worth generalising because negative results about binaries are usually
argued the weak way: entropy, section names, string searches, absence of known
signatures. Each of those is defeasible individually, and stacking defeasible
arguments does not produce an indefeasible one. A validated signature over the
current bytes does, because it converts ``we looked and found nothing'' into
``the file is provably identical to the one the publisher signed, so there is
nothing to find.''

The corollary is a cost argument: \textbf{the strong test was also the cheapest
one.} It was performed after the weak ones, which is the wrong order.

\subsection{Threats to validity}

\begin{itemize}
\item \textbf{Single system.} Every result comes from one laptop with one hybrid
      graphics configuration. The reassembly and forensic findings should
      transfer directly; the graphics findings are specific to this driver
      version, this translation layer, and this vendor pair.
\item \textbf{The root cause is localised, not explained.} We show that one build
      of the translation layer produces shader code that crashes a specific
      driver's compiler and that another does not. We did not identify the
      offending construct, and did not reduce it to a minimal reproducer. That
      would be the natural next step and would produce a filable driver bug.
\item \textbf{The unresolved contradiction of Section~\ref{sec:shaders}} is
      genuinely unresolved. The partial improvement from shader substitution has
      no confirmed mechanism.
\item \textbf{Timing figures} are wall-clock on a loaded desktop system, not a
      controlled benchmark.
\end{itemize}

\subsection{Related work}

Entropy analysis for packing detection follows Lyda and
Hamrock~\cite{lyda2007}; the windowing refinement and the known false-positive
behaviour are discussed at length in the packer-detection
literature~\cite{sokpacker}. Authenticode structure and validation semantics
are specified by Microsoft~\cite{authenticode} over PKCS\#7/CMS~\cite{rfc5652}
and the PE format~\cite{pecoff}. The bisection discipline in
Section~\ref{sec:lessons} is delta debugging~\cite{zeller2002}, and the general
approach to systematic debugging follows Zeller~\cite{zellerbook}. The
compatibility stack under test is Wine~\cite{wine} with Proton~\cite{proton},
vkd3d-proton~\cite{vkd3d}, and DXVK~\cite{dxvk}, targeting
Vulkan~\cite{vulkan} and SPIR-V~\cite{spirv}. The reassembly toolchain uses
LZ4~\cite{lz4}, VCDIFF~\cite{rfc3284}, HDiffPatch~\cite{hdiffpatch}, and Ogg
Vorbis~\cite{vorbis}; the installer format is Inno Setup~\cite{innosetup} and
the object-namespace behaviour is documented in Windows
Internals~\cite{ntobj}. Vendor dispatch is handled by libglvnd~\cite{glvnd}.

To our knowledge the specific combination reported here---a signature-based
negative determination of anti-tamper presence, a manual reconstruction of a
delta-compressed redistribution without its installer, and a cross-vendor shader
divergence causing a driver-side compiler fault under translation---has not been
described together in one place.

\section{Conclusion}
\label{sec:conclusion}

Four questions were asked and four were answered.

\textbf{Is it hostile?} No. Six independent lines of static evidence, all
negative, including a byte-for-byte comparison against an independently verified
package. The most alarming artefact in the distribution---a utility that rewrites
the system \verb|hosts| file---turned out to be the strongest single piece of
evidence \emph{for} the package's authenticity.

\textbf{Is commercial anti-tamper present?} No, and it never was in this build.
The executable's Authenticode digest validates over its current bytes, so the
file is exactly what the publisher signed. Protection of that class lives inside
the executable and cannot be removed without invalidating the signature.
Three corroborating instruments agree, including a windowed entropy profile with
zero windows anywhere near the threshold.

\textbf{Can it be reconstructed without its installer?} Yes---in 37 minutes, to a
working 54\,GB tree, with every one of ten stages verified in advance against a
manifest the format publishes for free. The artefact that blocked the project for
two working sessions was a three-line advertisement fetcher containing no build
logic at all, and it was blocked not by a codec but by a missing configuration
file at a hard-coded path.

\textbf{Will it run on both GPUs?} Yes. The integrated GPU worked early. The
discrete GPU produced a failure with no error message, no crash, and no log
output, which turned out to be nine unpainted modal dialogs each reporting that
the vendor's shader compiler had segmentation-faulted. The cause was neither the
game, nor the driver, nor the hardware, nor the permissions gate that was
suspected three separate times. It was the build of the Direct3D~12 translation
layer---the one component chosen at the start of the project and never once
treated as a variable. Replacing it took one attempt.

The recurring lesson is not about graphics. It is that in a stack of a dozen
substitutable layers, the expensive mistake is rarely a wrong hypothesis; it is a
component that was never a hypothesis. Twice in this project---the installer
helper that re-extracted itself, and the translation layer that was assumed
fixed---the answer was in a place that had been silently excluded from the search
space at the outset.

Two smaller results are worth carrying away independently. A validated code
signature is a stronger, cheaper, and more honest way to establish the absence of
anti-tamper protection than any heuristic stack, and it should be the first test
rather than the last. And a pipeline stage that publishes its own preconditions
gives you a correctness check you neither had to design nor have to trust
yourself to have designed correctly.

\section*{Appendix: Reproduction}
\label{sec:appendix}

\subsection*{A. Establishing anti-tamper absence on an arbitrary binary}

\begin{enumerate}
\item Extract the certificate directory from the PE file, parse the PKCS\#7
      structure, and read the embedded digest.
\item Recompute the Authenticode hash over the file as it exists, excluding the
      checksum field, the certificate table entry, and the certificate data.
      Compare. If they match, stop---the question is answered.
\item Otherwise, profile entropy in windows of 1\,MiB over the code section.
      Report the window count above 7.5 and state the threshold used. Never
      report a whole-file figure alone.
\item Enumerate sections and the import table. A full, compiler-shaped import
      table is difficult for a protection layer to preserve.
\item Identify what protection \emph{is} present. ``No anti-tamper'' is not ``no
      DRM.''
\end{enumerate}

\subsection*{B. Recovering a compiled Inno Setup script}

Search for the 12-byte identifier \verb|rDlPtS|; read nine little-endian 32-bit
words; seek to \textbf{offset0}; skip the 64-byte banner; read CRC, stored size,
and compression flag; then read \verb|[crc32(4) + up to 4096 bytes]| chunks until
the stored size is satisfied. The concatenated result is raw LZMA1 whose first
five bytes are the properties block. Do \emph{not} follow the \verb|zlb\x1a|
magic near the start of the overlay; that is the file data.

\subsection*{C. Reassembling the distribution}

\begin{enumerate}
\item Copy the decompressor's configuration file to the root of the emulated
      system drive.
\item Patch the compression helpers to drop the \verb|Global\| namespace prefix,
      using a live patcher that polls for newly extracted copies.
\item Extract all six containers to their \emph{correct} destinations
      (Table~\ref{tab:dest}); two of the six do not go where the others do.
\item Read the worklists from the smallest container for the commands, and the
      carved installer script for the ordering and working directories.
\item Before each directory patch, read its manifest and record the expected
      input count, input byte total, and output size. Verify after.
\item Parallelise the per-file worklists; the cost is process startup, not
      compression.
\end{enumerate}

\subsection*{D. Bisecting a translation-layer failure}

\begin{enumerate}
\item Enumerate windows inside the compatibility layer before assuming a hang.
      A modal dialog that is never painted is indistinguishable from a deadlock
      by every other means.
\item Read assertion text precisely. An NTSTATUS is not an API return code, and
      the difference decides whether you are debugging a rejection or a crash.
\item Dump shaders on both vendors and compare hashes \emph{and} bytes. Identical
      hashes with differing bytes means the layer is emitting vendor-specific
      code and your cross-vendor control is not a control.
\item Vary the translation-layer build before exhausting its configuration
      surface, not after.
\item Clone the distribution with hardlinks and replace libraries by
      delete-then-copy. Never edit in place. Give each variant its own prefix,
      and copy the prefix's sibling metadata, not just the prefix.
\end{enumerate}

\subsection*{E. Catalogue of silent failures encountered}

Each of the following reported success, or reported nothing, while doing
something other than what was intended.

\begin{table}[htbp]
\caption{Silent failures encountered in this project}
\label{tab:silent}
\centering
\footnotesize
\begin{tabularx}{\columnwidth}{@{}L{3.1cm}X@{}}
\toprule
\textbf{Construct} & \textbf{What it actually does} \\
\midrule
\verb|grep -c| with a fallback & prints \verb|0| \emph{and} exits non-zero, so
  the fallback appends a second zero; all later comparisons error and read false \\
Process match on full name & \verb|comm| truncates to 15 characters; the pattern
  never matches and the process reads as dead \\
Substring match on command line & matches the searching shell itself \\
Static patch of installer helpers & bypassed; the installer re-extracts fresh
  copies at runtime \\
Prefix seeded from \verb|pfx/| only & missing sibling metadata aborts the
  launcher before the game starts \\
Library search under \verb|lib64| & the libraries are under \verb|lib|; the
  search returns a confident false negative \\
GUID searched in printed order & the first three fields byte-swap in memory \\
Emptiness test on a \verb|/sys| file & always false; sysfs files report zero size
  regardless of content \\
Graphics layer's adapter report & deliberately reports AMD while on NVIDIA \\
Launcher's own notification & claimed NVIDIA while running on the integrated GPU \\
EGL vendor pin & replaces the vendor list rather than extending it, disabling an
  unrelated API for an unrelated vendor \\
Setgid helper raising only the effective GID & the shell drops it at startup \\
Archive tool's size counter & 32-bit signed; goes negative past 2\,GB \\
Debugger window text & truncated to 14 characters; labels look corrupted \\
Compositor geometry versus screenshots & two different coordinate systems
  differing by the scale factor \\
\bottomrule
\end{tabularx}
\end{table}

\subsection*{F. Artefacts}

The harness scripts (launcher, single-configuration trial, distribution variant
builder, three configuration sweeps, and a shader bisector), the archive
extraction harness and its freestanding 32-bit build recipe, the installer-script
carver, the helper patcher, and the reassembly drivers for both data chunks are
retained alongside a running findings log. Defect and decision records are
maintained in the project's own issue and decision registers.


\bibliographystyle{IEEEtran}
\begin{thebibliography}{99}

\bibitem{lyda2007}
R.~Lyda and J.~Hamrock, ``Using entropy analysis to find encrypted and packed
malware,'' \emph{IEEE Security \& Privacy}, vol.~5, no.~2, pp.~40--45, 2007.

\bibitem{sokpacker}
X.~Ugarte-Pedrero, D.~Balzarotti, I.~Santos, and P.~G. Bringas, ``SoK: Deep
packer inspection: A longitudinal study of the complexity of run-time packers,''
in \emph{Proc. IEEE Symp. Security and Privacy}, 2015, pp.~659--673.

\bibitem{authenticode}
Microsoft Corporation, ``Windows Authenticode portable executable signature
format,'' Technical Specification, 2008.

\bibitem{rfc5652}
R.~Housley, ``Cryptographic Message Syntax (CMS),'' RFC 5652, IETF, 2009.

\bibitem{pecoff}
Microsoft Corporation, ``PE format specification,'' Microsoft Learn, 2024.

\bibitem{zeller2002}
A.~Zeller and R.~Hildebrandt, ``Simplifying and isolating failure-inducing
input,'' \emph{IEEE Trans. Software Engineering}, vol.~28, no.~2, pp.~183--200,
2002.

\bibitem{zellerbook}
A.~Zeller, \emph{Why Programs Fail: A Guide to Systematic Debugging}, 2nd ed.
Morgan Kaufmann, 2009.

\bibitem{wine}
The Wine Project, ``Wine: a compatibility layer for running Windows applications
on POSIX-compliant operating systems,'' \url{https://www.winehq.org}.

\bibitem{proton}
Valve Corporation, ``Proton: compatibility tool for Steam Play,''
\url{https://github.com/ValveSoftware/Proton}.

\bibitem{vkd3d}
H.-K. Nienhuys \emph{et al.}, ``vkd3d-proton: Direct3D 12 to Vulkan translation
layer,'' \url{https://github.com/HansKristian-Work/vkd3d-proton}.

\bibitem{dxvk}
P.~Möller, ``DXVK: Vulkan-based implementation of D3D9, D3D10 and D3D11,''
\url{https://github.com/doitsujin/dxvk}.

\bibitem{spirv}
The Khronos Group, ``SPIR-V specification, version 1.6,'' 2023.

\bibitem{vulkan}
The Khronos Group, ``Vulkan 1.4 specification,'' 2024.

\bibitem{lz4}
Y.~Collet, ``LZ4: extremely fast compression algorithm,''
\url{https://github.com/lz4/lz4}.

\bibitem{rfc3284}
D.~Korn, J.~MacDonald, J.~Mogul, and K.~Vo, ``The VCDIFF generic differencing
and compression data format,'' RFC 3284, IETF, 2002.

\bibitem{hdiffpatch}
Housisong, ``HDiffPatch: a C\texttt{++} library and command-line tool for diff
and patch of binary files and directories,''
\url{https://github.com/sisong/HDiffPatch}.

\bibitem{vorbis}
Xiph.Org Foundation, ``Vorbis I specification,'' 2020.

\bibitem{glvnd}
NVIDIA Corporation and contributors, ``libglvnd: the GL vendor-neutral
dispatch library,'' \url{https://gitlab.freedesktop.org/glvnd/libglvnd}.

\bibitem{ntobj}
M.~Russinovich, D.~Solomon, and A.~Ionescu, \emph{Windows Internals, Part 2},
6th ed. Microsoft Press, 2012.

\bibitem{innosetup}
J.~Russell, ``Inno Setup,'' \url{https://jrsoftware.org/isinfo.php}.

\end{thebibliography}

\end{document}
